\documentclass[JEP,XML,SOM,Unicode,NoEqCountersInSection,published]{cedram}
\datereceived{2023-03-13}
\dateaccepted{2023-03-22}
\dateepreuves{2023-03-29}

\usepackage[scr=rsfs,cal=euler]{mathalfa}
\multlinegap0pt
\newenvironment{enumeratei}
{\bgroup\def\theenumi{\roman{enumi}}\def\theenumii{\arabic{enumii}}\begin{enumerate}}
{\end{enumerate}\egroup}

\def\mto{\mathchoice{\longmapsto}{\mapsto}{\mapsto}{\mapsto}}
\def\hto{\mathrel{\lhook\joinrel\to}}

\newtheorem{thmA}{Theorem}
\renewcommand{\thethmA}{\Alph{thmA}}

\usepackage{esint}

\newtheorem{theorem}{Theorem}[subsection]
\newtheorem{lemma}[theorem]{Lemma}
\newtheorem{prop}[theorem]{Proposition}
\newtheorem{conjecture}[theorem]{Conjecture}
\newtheorem{corollary}[theorem]{Corollary}
\theoremstyle{definition}
\newtheorem{example}[theorem]{Example}
\newtheorem{remark}[theorem]{Remark}
\newtheorem{defi}[theorem]{Definition}

\newcommand\bbc{\mathbb{C}}
\newcommand\bbd{\mathbb{D}}
\newcommand\bbp{\mathbb{P}}
\newcommand\bbr{\mathbb{R}}
\newcommand\bbs{\mathbb{S}}
\newcommand\cC{\mathcal{C}}
\newcommand\cE{\mathcal{E}}
\newcommand\cJ{\mathcal{J}}
\newcommand\cL{\mathcal{L}}
\newcommand\cM{\mathcal{M}}
\newcommand\cN{\mathcal{N}}
\newcommand\cO{\mathcal{O}}
\newcommand\cP{\mathcal{P}}
\newcommand\cR{\mathcal{R}}
\newcommand\cX{\mathcal{X}}
\newcommand\cY{\mathcal{Y}}
\newcommand\cZ{\mathcal{Z}}
\newcommand\K{\mathrm{K}}
\newcommand\R{\mathrm{R}}
\newcommand\mi{^{-1}}
\newcommand\an{^{\mathrm{an}}}
\newcommand\pdisc{\overline\bbd^*}
\newcommand\disc{\overline\bbd}

\DeclareMathOperator{\ord}{ord}
\DeclareMathOperator{\Specop}{Spec}
\newcommand\Spec[1]{\Specop \mathrm{#1}}
\newcommand\norm{\|{\cdot}\|}
\newcommand\PSH{\mathrm{PSH}}
\newcommand\MA{\,\mathrm{MA}}
\newcommand\refmetric{\phi_{\mathrm{ref}}}
\newcommand\na{^{\mathrm{NA}}}
\newcommand\ideal{\mathfrak{a}}

\datepublished{2023-03-31}
\begin{document}
\frontmatter
\title[The space of finite-energy metrics over a degeneration of complex manifolds]{The space of finite-energy metrics over\\ a degeneration of complex manifolds}

\author[\initial{R.} \lastname{Reboulet}]{\firstname{Rémi} \lastname{Reboulet}}
\address{Department of mathematical sciences, Chalmers University of Technology\\
412 96 Gothenburg, Sweden}
\email{rebouletremimath@gmail.com}
\urladdr{https://rebouletremimath.github.io/}

\begin{abstract}
Given a degeneration of projective complex manifolds $X\to \mathbb{D}^*$ with meromorphic singularities, and a relatively ample line bundle $L$ on $X$, we study spaces of plurisubharmonic metrics on $L$, with particular focus on (relative) finite-energy conditions. We endow the space $\hat {\mathcal{E}}^1(L)$ of relatively maximal, relative finite-energy metrics with a $d_1$-type distance given by the Lelong number at zero of the collection of fiberwise Darvas $d_1$-distances. We show that this metric structure is complete and geodesic. Seeing $X$ and $L$ as schemes $X_{\mathrm{K}}$, $L_{\mathrm{K}}$ over the discretely-valued field $\mathrm{K}=\mathbb{C}(\!(t)\!)$ of complex Laurent series, we show that the space $\mathcal{E}^1(L_{\mathrm{K}}^{\mathrm{an}})$ of non-Archimedean finite-energy metrics over $L_{\mathrm{K}}^{\mathrm{an}}$ embeds isometrically and geodesically into $\hat {\mathcal{E}}^1(L)$, and characterize its image. This generalizes previous work of Berman-Boucksom-Jonsson, treating the trivially-valued case.
\end{abstract}

\subjclass{32U05, 32Q15, 14E99}

\keywords{Berkovich spaces, complex manifolds, pluripotential theory, degenerations}

\altkeywords{Espaces de Berkovich, variétés complexes, théorie du pluripotentiel, dégénérescences}

\alttitle{L'espace des métriques d'énergie finie sur une dégénérescence de variétés complexes}

\begin{altabstract}
Étant donné une dégénérescence de variétés projectives complexes $X\to \mathbb{D}^*$ avec des singularités méromorphes, et un fibré en droites relativement ample $L$ sur $X$, nous étudions des espaces de métriques plurisousharmoniques sur $L$, avec une attention particulière aux conditions (rela\-tives) d'énergie finie. Nous munissons l'espace $\hat {\mathcal{E}}^1(L)$ des métriques relativement maximales d'énergie finie d'une distance de type $d_1$ donnée par le nombre de Lelong en $0$ de la famille des distances de Darvas $d_1$ fibre à fibre. Nous montrons que cette structure métrique est complète et géodésique. En considérant $X$ et $L$ comme des schémas $X_{\mathrm{K}}$, $L_{\mathrm{K}}$ sur le champ discrètement valué $\mathrm{K}=\mathbb{C}(\!(t)\!)$ des séries de Laurent complexes, nous montrons que l'espace $\mathcal{E}^1(L_{\mathrm{K}}^{\mathrm{an}})$ des métriques non archimédiennes d'énergie finie sur $L_{\mathrm{K}}^{\mathrm{an}}$ s'immerge isométriquement et géodésiquement dans $\hat {\mathcal{E}}^1(L)$, et nous caractérisons son image. Ceci généralise un travail précédent de Berman-Boucksom-Jonsson, traitant le cas trivialement valué.
\end{altabstract}

\maketitle
\tableofcontents
\mainmatter

\section*{Introduction}

\subsubsection*{Overview}

Given a finite-dimensional complex vector space $V$, the space of Hermitian norms on $V$ can be endowed with a metric space structure, by taking an $L^p$-norm of the vector of eigenvalues of the transition matrix between two such norms $\norm_0,\norm_1$; that is, the matrix whose entries are of the form $\langle e_i,e_j\rangle_1$, where $\langle \cdot,\cdot\rangle_1$ is the Hermitian product associated to $\norm_1$, and $(e_i)_i$ is an orthonormal basis for $\norm_0$. This makes the space of Hermitian norms into a \textit{negatively curved} metric space, implying that geodesic rays always exist; a natural question is then to understand the behaviour of such geodesic rays at infinity. In this case, there is a canonical isomorphism between equivalence classes of geodesic rays (where equivalence means that their distance remains bounded at infinity) of Hermitian norms, and \textit{non-Archimedean norms} on $V$, \ie nonnegative positive functions $\norm:V\to \bbr_{\geq 0}$ such that $\|\lambda\cdot v\|=\|v\|$ for $\lambda\in\bbc^\times$, $v\in V$, and satisfying the \textit{non-Archimedean} or \textit{ultrametric inequality} $\|v+w\|\leq \max(\|v\|,\|w\|)$, refining the usual triangle inequality. Furthermore, asymptotics of functionals along geodesic rays of Hermitian norms are captured by simpler functionals on the space of non-Archimedean norms, see \eg \cite[Th.\,1.6]{bouicm}.

A point of view on the present paper is that it aims at generalizing those results in a more geometric setting, where we no longer consider metrics on a vector space but rather on complex projective manifolds (we will consider the space of \textit{plurisubharmonic metrics}, which can be thought of as generalizing the notion of a Kähler metric); and where the time-variable is, in a sense, replaced by a \textit{holomorphic} time variable, allowing the complex structure of the manifold to vary nontrivially.

The study of rays of plurisubharmonic metrics has been, for the past decade, a~central component in successful approaches to important conjectures in complex geo\-metry. Given a polarized compact Kähler manifold $(X,L)$, we define a psh ray to be a psh metric on $L\times \pdisc$ over the trivial product $X\times\pdisc$, which is furthermore invariant under the usual action of $\bbs^1$. Setting the variable $t=-\log|z|$, we obtain a ``psh curve''
\[
[0,\infty)\ni t\mto \phi_t,
\]
and the study of the asymptotics of this curve (and of various quantities it can be evaluated against) plays a crucial role in many questions. One can picture such a curve as a way to deform $(X,L,\phi_0)$. A particular striking application has been a proof of the Kähler-Einstein case of the Yau-Tian-Donaldson conjecture due to Berman-Boucksom-Jonsson (\cite{bbj}), and further advances in the general cscK case by Li (\cite{liytd}), both of which are major inspirations for this current work. In parallel, Darvas-Lu have started a very interesting investigation of the metric properties of the space of such rays (\cite{darlurays}), which has found many interesting developments.

One side of the story, very apparent in the aforementioned proof of the Yau-Tian-Donaldson conjecture, is the introduction of valuative (non-Archimedean) techniques to study asymptotics in Kähler geometry. Namely, to a psh ray $t\mto \phi_t$ as above, one can associate a non-Archimedean function $\phi\na$, defined on the Berkovich \hbox{analytification} $X\an$ of $X$ with respect to the trivial absolute value on $\bbc$. Without going into too much detail, $X\an$ can be essentially described as the space of valuations on~$\bbc(X)$, which is then compactified using semivaluations, so that $\phi\na$ can be thought of as capturing the asymptotic singularities of $\phi$ along all possible birational~models of~$X$. A beautiful result is that a class of distinguished geodesic rays (the ones of interest for K-stability) are in one-to-one correspondence with non-Archimedean functions of finite energy on $X\an$. Thus, the space of non-Archimedean functions can be thought of as a boundary at infinity for the space of Archimedean plurisubharmonic functions.

In the present work, we generalize such results to the wider setting of psh metrics on arbitrary families of projective manifolds over the punctured disc. Already in the isotrivial case without $\bbs^1$-action, these spaces have garnered interest (see \eg \cite{donaldsondiscs}, \cite{rwndiscs}), but our results also hold for arbitrary degenerations with ``meromorphic'' singularities. This setting encompasses \eg families of Kähler-Einstein metrics, which have been studied for well over a decade (see \cite{sch08}, \cite{tsujidynamical}, \cite{cgp2}...), but also works related to the Kähler-Ricci flow, to collapsing families of Calabi-Yau manifolds... We refer the reader to Section \ref{subsect_misc} for some discussion and new questions in various directions. A main interest of allowing such general families lies, for example, in Donaldson's example (\cite[\S 5.1]{donbstability}) of a special fiber not accessible via test configurations.

\subsubsection*{The complex side of the story: the space of degenerating metrics}

We will be working on meromorphic degenerations (or simply degenerations) $\pi:X\to\pdisc$, which are defined as follows: $\pi:X\to\pdisc$ is a proper holomorphic submersion endowed with a relatively ample line bundle $L\to X$ (in particular, the fibers of $\pi$ are projective), which we furthermore assume to be the restriction of a normal complex analytic space~$\cX$ fibered over $\disc$. We shall say that such a space~$\cX$ is a \textit{model} of $X$, and call its fiber~$\cX_0$ over zero its \textit{central} or \textit{special fiber}. If $\cL$ is a relatively ample line bundle on a model~$\cX$, isomorphic to $L$ away from $\cX_0$ and compatible with the identification $\cX-\cX_0\simeq X$, we say that $(\cX,\cL)$ is a \textit{model of $(X,L)$}. We will be interested in spaces of psh metrics on line bundles $L$ over a (meromorphic) degeneration $X$.

Fixing a fiber $X_z$ for the moment, Darvas introduces in \cite{darv} a $L^1$-type metric structure on the space of finite-energy metrics $\cE^1(X_z,L_z)$, where the distance between two metrics $\phi_{0,z},\phi_{1,z}$ is given by
\[
d_{1,z}(\phi_{0,z},\phi_{1,z})=E(\phi_{0,z})-E(\phi_{1,z})+2E(P(\phi_{0,z},\phi_{1,z})).
\]
The term $E$ is the Monge-Ampère (or Aubin-Yau) energy of a psh metric, which is a normalized primitive of the Monge-Ampère operator
\[
\phi_z\mto (dd^c\phi_z)^d
\]
(where $d=\dim X_z$), for $\phi_z$ smooth, and which is extended to general finite-energy metrics via decreasing limits, as in \cite{bbgz}.
On the other hand, the $P$ term is the ``rooftop'' envelope
\[
P(\phi_{0,z},\phi_{1,z})=\sup\{\phi_z\in\PSH(X_z,L_z),\,\phi_z\leq\phi_{0,z},\phi_{1,z}\},
\]
which generalizes the convex envelope of the minimum of two convex functions. If the two metrics in the left-hand side have finite energy, then their rooftop envelope can also be shown to have finite energy, so that all the terms in the above expression are finite.

We wish to define a distance on the space of metrics on $L$ (with fiberwise finite energy) via taking the Lelong number of the collection of fiberwise distances above. This is to be understood as a generalized (signed) Lelong number as follows: given a subharmonic function $f$ on the punctured disc such that $g:=f+a\log|z|$ is bounded above for some $a$, then $g$ is also subharmonic on the whole disc with finite Lelong number at zero, and we define $\nu_0(f):=\nu_0(g)-a$.
Assuming $z\mto d_1(\phi_{0,z},\phi_{1,z})$ to satisfy this claim, we then define the distance $\hat d_1(\phi_0,\phi_1)$ to be its generalized Lelong number. As the reader can see, we have made two assumptions: that the distance between $\phi_{0,z}$ and $\phi_{1,z}$ has at most logarithmic growth, and that this distance function is subharmonic on the disc.

The first assumption is implied by a natural condition on such metrics $\phi$, which we abusively also call \textit{logarithmic growth}, and can equivalently be described as requiring that $\phi$ extends plurisubharmonically over the central fiber of some analytic model of $X$. This is a very general condition and really the necessary minimum to ask for; therefore, we will denote $\PSH(L)$ the space of psh metrics of logarithmic growth on~$L$. In the case of rays, this corresponds to the usual linear growth condition (as in \cite{bbj}).

Regarding the subharmonicity statement, if one looks at the $\bbs^1$-invariant case, similar statements (see \cite{bdl}) require that rays be \textit{geodesic}. This suggests that in our setting, we must look at maximality conditions coming from the base. In our setting, this is understood in a ``relative'' pluripotential theory sense, \ie a metric~$\phi$ with logarithmic growth on $L$ is relatively maximal if for any relatively compact open set $U$ on the base $\pdisc$, $\phi|_{\pi\mi(U)}$ is larger than all psh metrics on $L|_U$ bounded above by $\phi$ on $\partial U$. The space of fiberwise finite-energy, logarithmic growth metrics on~$L$ is denoted by $\cE^1(L)$, the subspace of $\cE^1(L)$ corresponding with metrics which are furthermore \textit{relatively maximal} is then denoted by $\hat\cE^1(L)$ Note that, although there appears to be many assumptions, this is exactly the generalization of the setting considered in \cite{bbj} and \cite{darlurays}; more precisely, assuming $\bbs^1$-invariance, our space $\cE^1(L)$ corresponds to the space of (non-necessarily geodesic) psh rays, while $\hat\cE^1(L)$ corresponds with the space $\cR^1(L)$ of \cite{darlurays}, and our distance $\hat d_1$ corresponds with the distance
\[
d_1^\cR(\phi_0,\phi_1):=\lim_{t\to\infty}t\mi d_1(\phi_{0,t},\phi_{1,t}).
\]
We then show the following:

\begin{thmA}[\ref{thm_1mainthmsepdiscs}, \ref{thm_1completeness}]\label{th:A}
The space $(\hat\cE^1(L),\hat d_1)$ is a complete metric space.
\end{thmA}

In a nutshell, this result means that we obtain a genuine ``space of families of metrics'' in which one can use the tools of metric and variational analysis. We obtain a natural notion of distance between such families, which allows us to say how ``close'' they are to each other in a geometrically meaningful way, and also helps in formulating fundamental notions such as compactness.

\subsubsection*{Going non-Archimedean}

As explained for example in \cite{favre}, one can interpret a projective (meromorphic) degeneration $X$ as a projective variety over the field $\K=\bbc(\!(t)\!)$. This is a discretely-valued field, \ie the possible values of its usual valuation form a discrete additive subgroup of $\bbr$. As in the trivially valued case discussed at the very beginning, we define a metric $\phi\na$ on the Berkovich analytification $L_\K\an$ of $L$ seen as a $\K$-line bundle, associated to a metric $\phi$ in $\PSH(L)$, which captures the algebro-geometric information of the singularities of $\phi$ via generic Lelong numbers along vertical divisors in models of $X$. We prove that the metric $\phi\na$ is plurisubharmonic in the non-Archimedean sense (see Section \ref{sect_nappt}), which roughly means that it satisfies the statement of Demailly's regularization theorem, \ie can be approximated by a decreasing net of Fubini-Study metrics.

\subsubsection*{Energies and Deligne pairings}

More structure arises when one considers finite-energy metrics. For this, it will be more convenient to express complex and non-Archimedean Monge-Ampère energies as metrized Deligne pairings (see Section \ref{sect_deligne}). Briefly, one can associate to a flat projective morphism of complex manifolds
\[
\pi:X\to Y
\]
of relative dimension $d$, and to $d+1$ pairs $(L_i,\phi_i)$ of relatively ample line bundles $L_i$ over $X$ together with continuous psh metrics $\phi_i$, a line bundle
\[
\langle L_0,\dots,L_d\rangle_{X/Y}
\]
together with a metric
\[
\langle \phi_0,\dots,\phi_d\rangle_{X/Y}
\]
in a multi-additive and symmetric fashion. Furthermore, this construction satisfies
\begin{itemize}
\item a change of metric formula: given another continuous psh metric $\phi_0'$ on $L_0$,
\begin{equation}
\langle \phi_0,\dots,\phi_d\rangle_{X/Y} - \langle \phi_0',\dots,\phi_d\rangle_{X/Y} = \pi_*\left((\phi_0-\phi_0')(dd^c\phi_1\wedge\dots\wedge dd^c\phi_d)\right);
\end{equation}
\item and a pushforward formula:
\begin{equation}
dd^c\langle \phi_0,\dots,\phi_d\rangle_{X/Y}=\pi_*(dd^c\phi_0\wedge\dots\wedge dd^c\phi_d).
\end{equation}
\end{itemize}
If $Y$ is a point, we omit the subscript $\cdot_{X/Y}$, and upon contemplating the change of metric formula, one can see that the (relative) Monge-Ampère energy between two metrics $\phi_0$ and $\phi_1$ may be written as
\[
\langle \phi_0^{d+1}\rangle - \langle \phi_1^{d+1}\rangle,
\]
suggesting that the Monge-Ampère energy could be seen as a genuine metric $\langle \phi^{d+1}\rangle$.

In the present article, we also extend the Deligne pairing construction in the complex setting, from continuous psh metrics to fiberwise finite-energy metrics. We direct the interested reader to Section \ref{sect_deligne}.

\subsubsection*{Perfecting the complex/non-Archimedean correspondence}

In the space of non-Archimedean psh metrics, one can also define the subclass of finite-energy metrics $\cE^1(L\an)$, by defining for example the Monge-Ampère energy $E\na$ as a non-Archimedean Deligne pairing (using the machinery of \cite{be20}). (Other approaches, which all give the same functional, include the intersection theory-based formalism of Gubler (\cite{gublerinventionestropical}, and the locally tropical approach of Chambert-Loir and Ducros (\cite{cld}) based on the superforms of Lagerberg (\cite{lagerberg}).) A natural question is now whether the Monge-Ampère energy is equivariant under the ``Archimedean to non-Archimedean'' map. We can write this statement as
\begin{equation}\label{eq_ext}E(\phi)\na=E\na(\phi\na)
\end{equation}
(seeing Monge-Ampère energies as Deligne pairings, as previously), and can understand it as saying that the generalized Lelong number at zero of the Monge-Ampère energy along a metric $\phi$ in $\hat\cE^1(L)$ coincides with the non-Archimedean Monge-Ampère energy of $\phi\na$, in a sense made precise in Remark \ref{rem_namap}. It turns out that, in general, we only obtain an inequality
\[
E(\phi)\na\leq E\na(\phi\na),
\]
as in Theorem \ref{thm_uniquenesstropmaxdisc}. As pointed out by the anonymous referee, this inequality has a local analogue using the valuative transform of Boucksom-Favre-Jonsson. Let $\phi$ be a germ of a psh function near $0$ in $\bbc^d$, and $\hat\phi$ is its valuative transform (in the sense of \cite[\S 5.2]{bfjlocal}) on the space of normalized valuations centered at zero. Then, the inequality above corresponds to the fact that
\[
(dd^c \phi)^d(\{0\})\geq \MA\na(\hat\phi),
\]
where $\MA\na$ is then the Monge-Ampère measure on this valuation space in the sense of \cite[\S 4.2]{bfjlocal}.

It becomes natural to take our interest to the class of metrics in $\hat\cE^1(L)$ for which equality holds. We define them as ``hybrid maximal'' metrics, which correspond to the maximal psh geodesic rays of \cite{bbj}. We give equivalent characterizations of such metrics, in particular using more ``complex pluripotential''-theoretic notions such as an extremal characterization. Denoting the space of hybrid maximal metrics by $\hat\cE^1_{\mathrm{hyb}}(L)\subset \hat\cE^1(L)$, we then prove that it can be completely realized by the space $\cE^1(L\an)$. The following theorem generalizes results of \cite[\S 6]{bbj}, \cite[\S 4]{liytd}.

\begin{thmA}[combining Theorems \ref{thm_isometry}, \ref{thm_uniquenesstropmaxdisc} \& \ref{prop_functionals}]\label{th:B}
Let $(X,L)$ be a meromorphic degeneration, and $(X\an,L\an)$ be the Berkovich analytifications of their associated varieties over the field $\bbc(\!(t)\!)$. We then have the following:
\begin{itemize}
\item there is an isometric embedding of $(\cE^1(L\an),d_1\na)$ into $(\hat\cE^1(L),\hat d_1)$ with image $\hat\cE^1_{\mathrm{hyb}}(L)$; this image is characterized as the class of metrics for which \eqref{eq_ext} holds;
\item in particular, $(\cE^1(L\an),d_1\na)$ and $\hat\cE^1_{\mathrm{hyb}}(L)$ are complete, geodesic metric spaces;
\item there is a general ``plurifunctional extension'' property in this space, as follows: suppose given $d+1$ relatively ample line bundles $L_i$ on $X$. Then, for any $(d+1)$-uple of metrics $\phi_i\in\hat\cE^1_{\mathrm{hyb}}(L_i)$, we have
\[
(\langle\phi_0,\dots,\phi_d\rangle_{X/\pdisc})\na=\langle\phi\na_0,\dots,\phi\na_d\rangle.
\]
\end{itemize}
\end{thmA}

This result states that the more important degenerations of psh metrics can be studied using purely non-Archimedean techniques. Conversely, we may also deduce properties of non-Archimedean metrics from their complex counterparts; non-Archimedean metrics are a priori quite mysterious objects, and Theorem~\ref{th:B} says that we can completely realize them using more familiar complex objects. This dictionary, combined with \eg the third point, is very powerful: in the $\bbs^1$-invariant case, testing slopes at infinity of certain functionals along a geodesic ray allows one to check for existence of the fabled constant scalar curvature Kähler metrics, and parts of such functionals are of the type given in point three. This result then states that one needs only to check these parts of the functionals on the ``easier'', more algebraic non-Archimedean side. We therefore expect our result to help in the study of similar problems \textit{in families}, such as existence of families of metrics satisfying Kähler-Einstein-type or cscK-type equations, as we discuss briefly in Section \ref{subsect_misc}.

In a different direction, the second point of Theorem~\ref{th:B} shows that the space of non-Archimedean functions has good topological properties; on the other side of the dictionary, the geodesicity statement gives a canonical way to transform a family of metric into another, and allows one to introduce a notion of \textit{convexity} directly in the space $\hat\cE^1_\mathrm{hyb}(L)$. In particular, one can now understand what it means for a functional to be (geodesically) convex in this space, and to speak of \eg their minimizers, which likely would have great geometric importance.

\subsubsection*{Organization of the paper}

In Section \ref{sect_sect1}, we study spaces of psh metrics with relative finite energy in the general relative case. In Section \ref{sect_sect15}, we specialize to degenerations over the punctured disc, and study our space $\hat\cE^1(L)$, building up to Theorem~\ref{th:A}. In~Section \ref{sect_sect2}, we construct the ``Archimedean to non-Archimedean'' map, after recalling some notions of non-Archimedean pluripotential theory. In~Section~\ref{sect_sect3}, we study the non-Archimedean limits of energy functionals, then hybrid maximal metrics, proving Theorem~\ref{th:B}. We also study the trivially-valued case in Section \ref{subsect_36}, and discuss other applications regarding families of Kähler-Einstein metrics.

\subsubsection*{Acknowledgements} The author thanks his advisors, Sébastien Boucksom and Catriona Maclean. He thanks Robert Berman, Bo Berndtsson, Tamás Darvas, Vincent Guedj, Léonard Pille-Schneider, and Steve Zelditch for various discussions. He also thanks the anonymous referee for many valuable remarks on the paper, in particular regarding clarity of exposition and for pointing out unclear points in certain proofs.

\section{Relative finite-energy spaces}\label{sect_sect1}

\subsection{Reminders on finite-energy spaces}

We begin with some reminders concerning $d_1$-structures on spaces of finite-energy metrics in the classical setting. We thus consider a fixed \textit{compact} Kähler manifold $X$, with $\dim X =: d$, endowed with an ample line bundle $L$ (so that $X$ is projective). Throughout this article, we will use additive conventions for tensor products of line bundles, meaning that if $L$ and $M$ are line bundles and $k$ a positive real integer, then $kL-M$ denotes $L^{\otimes k}\otimes M^{-1}$. Likewise, if $\phi$ and $\psi$ are metrics on $L$ and $M$, then we shall write $k\phi-\psi$ to denote the induced metric on $kL-M$.

Recall that a singular metric $\phi$ on $L$ is a Hermitian metric given in a trivialization $L|_U\simeq U\times \bbc$ by a function of the form $e^{-\phi_U}$ with $\phi_U\in L^1_{\mathrm{loc}}(U)$. We say that a singular metric $\phi$ is \textit{plurisubharmonic} or \textit{psh} if the current locally given by $i\partial\bar\partial \phi_U$ is positive. We define $\PSH(L)$ to be the set of psh metrics on $L$.

Consider two metrics $\phi_0$, $\phi_1\in C^0\cap \PSH(L)$. Their relative Monge-Ampère energy is the quantity
\[
E(\phi_0,\phi_1)=\frac{1}{(c_1(L)^d)(d+1)}\sum_{i=0}^{d}\int_X (\phi_0-\phi_1)\,(dd^c\phi_0)^i\wedge(dd^c\phi_1)^{d-i}.
\]
Note that we have a cocycle identity
\[
E(\phi_0,\phi_1)=E(\phi_0,\phi')+E(\phi',\phi_1)
\]
for any other continuous psh metric $\phi'$. Having fixed a continuous psh metric $\refmetric$ on the right, $E(\phi):=E(\phi,\refmetric)$ can be seen as an operator on $C^0\cap\PSH(L)$. This operator is also a primitive of the Monge-Ampère operator $\MA:\phi\mto (dd^c\phi)^d$ in the sense that
\[
\left.\frac{d}{dt}\right|_{t=0}E(\phi+tf)=\int_X f\,\MA(\phi)
\]
for $f\in C^0(X)$. It admits a (possibly infinite) extension to $\PSH(L)$ via
\[
E(\phi)=\lim_{k\to\infty}E(\phi_k),
\]
where $\phi_k$ is a net of continuous psh metrics decreasing to $\phi$, which always exists by \cite{demcurrents} (see also \cite{blockolo} for a simpler and more recent proof). The space of finite-energy metrics is the space
\[
\cE^1(L)=\{\phi\in \PSH(L),\,E(\phi)\text{ is finite}\}.
\]
By the cocycle identity, this space does not depend on the choice of a reference metric. From the work of Darvas, we know this space to admit a $d_1$-type complete metric space structure via
\[
d_1(\phi_0,\phi_1)=E(\phi_0)+E(\phi_1)-2E(P(\phi_0,\phi_1)),
\]
where $P(\phi_0,\phi_1)$ is the envelope
\[
P(\phi_0,\phi_1)=\sup\,\{\phi\in \PSH(L),\,\phi\leq\min(\phi_0,\phi_1)).
\]
It will be more practical to use a different expression of the Monge-Ampère energy, as a difference of absolute Deligne pairings.

\Subsection{Relative finite-energy metrics and extended Deligne pairings}\label{sect_deligne}

The Deligne pairing has a long history, starting from Deligne's original article, treating the case of relative dimension $1$ (\cite{deligne}), further generalized by Elkik in \cite{elkik1}, \cite{elkik2}. Its use to formulate functionals arising in complex geometry has been popularized via \cite{phongrosssturm}, and recently, Deligne pairings have also been shown to be of great use in non-Archimedean geometry (\cite{bhjasymptotics}, \cite{be20}, \cite[Rem.\,6]{phongrosssturm}), see also \cite{yz}. The non-Archimedean case over a point has been thoroughly developed in \cite{be20}.

We consider a proper holomorphic submersion between complex manifolds~\hbox{$\pi{:}\,X\!\to\! Y$} of relative dimension $d$. Pick $d+1$ pairs $(L_i,\phi_i)$, where $L_i$ is a line bundle over $X$, and $\phi_i$ is a continuous metric on $L_i$. To this data, one associates a line bundle over $Y$, $\langle L_0,\dots,L_d\rangle_{X/Y},$ together with a metric $\langle \phi_0,\dots,\phi_d\rangle_{X/Y}$ in a way that is multi-additive, symmetric; the construction furthermore commutes with base change (in particular, is stable upon restriction to an open set on the base), and satisfies
\begin{itemize}
\item the change of metric formula: given another continuous metric $\phi_0'$ on $L_0$, we have
\begin{equation}\label{eq_changeofmetricformula}
\langle \phi_0,\dots,\phi_d\rangle_{X/Y} - \langle \phi_0',\dots,\phi_d\rangle_{X/Y} = \pi_*\left((\phi_0-\phi_0')(dd^c\phi_1\wedge\dots\wedge dd^c\phi_d)\right)
\end{equation}
(see \cite[Th.\,I.1.1(d)]{elkik2});
\item the curvature formula
\begin{equation}\label{eq_ddc}
dd^c\langle \phi_0,\dots,\phi_d\rangle_{X/Y}=\pi_*(dd^c\phi_0\wedge\dots\wedge dd^c\phi_d)
\end{equation}
(see \cite[Th.\,I.1.1(d)]{elkik2}).
\end{itemize}
The last formula shows that the metric $\langle \phi_0,\dots,\phi_d\rangle_{X/Y}$ is positive if all the $\phi_i$ are psh. We also remark that, by multi-additivity, the Deligne pairing of $\cO_X$ with any other line bundles yields the trivial line bundle on the base, so that any metric on a Deligne pairing of the form $\langle \cO_X,L_1,\dots,L_d\rangle_{X/Y}$ can naturally be identified with a function on $Y$ (since metrics on the trivial line bundle are simply functions).

Assume for the moment that $Y$ is a point. In that case, Deligne pairings can be seen as complex lines together with a Hermitian norm. In this setting, we will omit the subscript $\cdot_{X/Y}$. Using the change of metric formula, one can see the relative Monge-Ampère energy between two continuous psh metrics on a fixed line bundle $L$ over $X$ as a difference of Deligne pairings:
\[
(d+1)E(\phi_0,\phi_1)=\langle \phi_0^{d+1}\rangle-\langle \phi_1^{d+1}\rangle,
\]
which suggests that the Monge-Ampère energy can be seen intrinsically as a genuine (Hermitian) metric $(d+1)E(\phi)=\langle \phi^{d+1}\rangle$ on the line $\langle L^{d+1}\rangle$.

We now return to arbitrary $Y$. Our goal is to extend the Deligne pairing construction to the class of metrics of finite-energy, in a way that the change of metric and curvature formulas still hold. To that end, we first introduce a new class of metrics allowing us to make sense of such formulas.
\begin{defi}Let $\pi:X\to Y$ be a proper holomorphic submersion between complex manifolds of relative dimension $d$. Let $L$ be a relatively ample line bundle on $X$. We define the class of relative finite-energy metrics $\cE^1_{X/Y}(L)$ to be the class of plurisubharmonic metrics $\phi$ on $L$ such that, for all $y\in Y$, $\phi_y\in\cE^1(L_y)$. Here, $L_y$ is the restriction of $L$ to the fiber $\pi\mi(y)$.
\end{defi}

Since we have required plurisubharmonicity on all of $L$, it follows that any metric in $\cE^1_{X/Y}(L)$ can be approximated by a decreasing net of continuous psh metrics on $L$. In particular, such metrics admit Deligne pairings.

\begin{theorem}\label{thm_eonepairing}Let $\pi:X\to Y$ be a proper holomorphic submersion between complex manifolds of relative dimension $d$, and let $(L_i)_{i=0}^d$ be a collection of $d+1$ relatively ample line bundles on $X$. There exists a unique extension of the Deligne pairing construction to metrics in $\cE^1(L_i)_{X/Y}$, giving a finite-valued metric on the Deligne pairing $\langle L_0,\dots,L_d\rangle_{X/Y}$ over $Y$, which is multilinear, symmetric, stable upon restriction to a smaller open set on the base, and such that the change of metric formula \eqref{eq_changeofmetricformula} holds.
\end{theorem}
Note that the curvature formula a priori will not hold in general, due to a lack of control of the right-hand side of \eqref{eq_ddc} along arbitrary decreasing limits.

Along the way, we will need the following lemma regarding finiteness of products of absolute finite-energy classes. This is a very standard result, which follows from \eg the arguments of \cite[Th.\,5.8]{bjtrivval21}, therefore we leave the details to the interested reader.
\begin{lemma}\label{lem_finitepairings}Let $X$ be a compact Kähler manifold of dimension $d$, and let $(L_i)$ be a collection of $d+1$ ample line bundles on $X$. Fix, for all $i=0,\dots,d$, a metric $\phi_i\in\cE^1(L_i)$, and a continuous metric $\phi'_0\in\cE^1(L_0)$. Then, the integral
\[
\int_X (\phi_0-\phi'_0)\,dd^c\phi_1\wedge\dots\wedge dd^c\phi_d
\]
is finite.
\end{lemma}
\begin{proof}[Proof of Theorem \ref{thm_eonepairing}.]
We first restrict to an open set $U$ on the base $Y$, so that we may apply Demailly regularization on $\pi\mi(U)$. Fix for each $i$ a metric $\phi_i\in\cE^1_{X/U}(L_i)$, and let $k\mto \phi_i^k$ be a sequence of continuous psh metrics on $L_i$ decreasing to $\phi_i$. We~claim that the sequence
\[
k\mto \langle \phi_0^k,\dots,\phi_d^k\rangle_{X/U}
\]
decreases to a finite-valued, finite-energy metric on $\langle L_0,\dots,L_d\rangle_{X/U}$, independent of the choices of approximating sequences, which defines our construction restricted~to~$U$. Assuming this convergence to hold, one sees that this construction is multilinear, symmetric, satisfies the change of metric formula. Uniqueness follows from the change of metric formula, which itself shows that the construction glues well over $X$.

That it would define a finite-valued metric on $U$ follows from Lemma \ref{lem_finitepairings} above, so that all that is left in order to prove the theorem is that the limit in question is decreasing. By the relative finite energy hypothesis, this limit will be finite and hence define a finite-energy metric on the pairing over $Y$. We proceed by induction on the number $n$ of indices $i\in\{0,\dots,d+1\}$ such that $\phi_i$ belongs to $\cE^1_{X/U}(L_i)-C^0\cap\PSH(L_i)$. In the case $n=0$, all metrics are continuous psh and this is the classical Deligne pairing, so that we have nothing more to prove.

Assume thus that the assertion holds for some $d+1>n>0$. Assume the metrics $\phi_i$, $i=1,\dots,d+1-n$ to belong to $C^0\cap\PSH(L_i)$, and the $n+1$ other metrics $\phi_i$ to belong strictly to $\cE^1_{X/U}(L_i)$, $i=0$ or $i=d+2-n,\dots,d$. (We can do this without loss of generality, by symmetry and up to reordering the indices.) We approximate $\phi_0$ and the $(\phi_i)_{i=d+2-n}^{d}$ by sequences $k\mto \phi_i^k$ of continuous psh metrics. For a fixed $\ell\in\mathbb{N}^*$ and by the induction assumption, the sequence
\[
k\mto \langle \phi_0^\ell,\phi_1,\dots,\phi_{d+1-n},\phi^k_{d+2-n},\dots,\phi^k_{d}\rangle_{X/Y}
\]
is decreasing and converges to a limit
$\langle \phi_0^\ell,\phi_1,\dots,\phi_{d+1-n},\phi_{d+2-n},\dots,\phi_{d}\rangle_{X/U}$. This limit satisfies, for any fixed metric $\phi'_0\in C^0\cap\PSH(L_0)$ the formula
\begin{multline*}
\langle \phi_0^\ell,\phi_1,\dots,\phi_{d+1-n},\phi_{d+2-n},\dots,\phi_{d}\rangle_{X/U}-\langle \phi'_0,\phi_1,\dots,\phi_{d+1-n},\phi_{d+2-n},\dots,\phi_{d}\rangle_{X/U}\\
=\int_{X/U}(\phi_0^\ell-\phi'_0)\,dd^c\phi_1\wedge\dots\wedge dd^c\phi_{d+1-n}\wedge dd^c\phi_{d+2-n}\wedge\dots\wedge dd^c\phi_{d}.
\end{multline*}
Now, this expression yields a decreasing net as $\ell$ increases, and its limit is finite. In particular, it can be seen to be the decreasing limit of
\[
k\mto \langle \phi_0^\ell,\phi_1,\dots,\phi_{d+1-n},\phi^k_{d+2-n},\dots,\phi^k_{d}\rangle_{X/U},
\]
which proves our desired statement by induction.
\end{proof}

\Subsection{Relatively maximal metrics}\label{sect_relmax}

\begin{defi}
Let $\pi:X\to Y$ be a holomorphic submersion with projective fibers, where $Y$ is a possibly open complex manifold. Let $L$ be a relatively ample line bundle on $X$. We say that a metric $\phi$ on $L$ is relatively maximal if for any relatively compact open subset $U$ of $Y$, for any relatively compact open subset~$V$ of $\pi\mi(U)$, and for any psh metric $\psi$ on the restriction of $L$ to $\pi\mi(V)$ such that $\limsup (\psi(z)-\phi(z))\leq 0$ as $z$ approaches the boundary of $\pi\mi(V)$, then
\[
\psi(z)\leq\phi(z)
\]
for all $z$ in $\pi\mi(V)$.
\end{defi}
Note that this means that, for all open sets $\pi\mi(U)$ as above, the restriction of $\phi$ to $\pi\mi(U)$ is \textit{maximal} in the sense of Sadullaev (\ie on such open sets, $\phi$ satisfies the usual definition of maximality as can be found in \eg \cite{klimek}).

\begin{remark}One sees from this definition that a decreasing limit of relatively maximal psh metrics is also relatively maximal.
\end{remark}

\begin{remark}
Let $M$ be a projective manifold together with an ample line bundle $L_M$. Let $[0,\infty)\ni t \mto \phi_t$ be a psh ray of psh metrics on $L_M$. We can identify $\phi$ with a $\bbs^1$-invariant psh metric $\tilde\phi$ on the product $L\times\pdisc$, by setting $\tilde\phi_z:=\phi_{-\log|z|}$ for $z\in\mathbb{D}^*$. Then, $\tilde\phi$ is relatively maximal in our sense if and only if $\phi$ is ``geodesic'' in the sense of \cite{bbj}.
\end{remark}

A nice way to generate relatively maximal metrics is via Perron-Bremmermann envelopes, as we now prove. We extend our setting slightly, to allow for singular fibers, which will be useful later on. We state our result in maximal generality, but the case to keep in mind is that of a holomorphic submersion over the punctured disc with a singular fiber over zero.
\begin{theorem}\label{thm_existencemaximal}
Let $\pi:X\to Y$ be a holomorphic projective surjective morphism. Let $\Omega$ be a relatively compact open subset of $Y$ with smooth boundary, such that $\pi$ is a submersion above (hence near) $\partial\Omega$. Let $L$ be a $\pi$-ample line bundle on $X$. Let $\phi$ be a continuous collection of fiberwise psh metrics on $\pi\mi(\partial\Omega)$. We then have that:
\begin{enumerate}
\item if there exists a continuous psh extension of $\phi$ to all of $\pi\mi(\Omega)$, then there exists a (unique) relatively maximal continuous psh extension of $\phi$ to all of $\pi\mi(\Omega)$;
\item if $\Omega$ is defined as $\{\rho < 0\}$, where $\rho$ is a smooth strictly psh function on $Y$, such that $\nabla \rho \neq 0$ whenever $\rho=0$, then a continuous psh extension as above exists.
\end{enumerate}
\end{theorem}
\begin{remark}An open subset that satisfies the second point above is called a \textit{hyperconvex} open subset. In particular, $\bbd$ and annuli centered at zero are such open sets. The proof of the first point follows some ideas dating back to the work of Bedford-Taylor (\cite{btdirichlet}), see \eg \cite[Prop.\,6.3]{bbgz}.
\end{remark}

\begin{proof}[Proof of Theorem \ref{thm_existencemaximal}.] The hypotheses in the theorem give that $\pi\mi(\bar \Omega)$ is an analytic space with boundary, which we denote $\bar M:=\pi\mi(\bar \Omega)$, and whose boundary is $\pi\mi(\partial \Omega)$, which we denote $\partial M:=\pi\mi(\partial\bar \Omega)$. We will finally write $M:=\pi\mi(\Omega)$.

\subsubsection*{Proof of the first point: existence of a continuous relatively maximal metric, assuming exis\-tence of a subsolution}
We assume that there exists a (weak) subsolution to the problem, \ie that there exists a metric $\psi\in C^0\cap \PSH(L|_M)$ which is equal to $\phi$ on $L|_{\partial M}$. (We will show that this is true in the second part of the proof.) Under this assumption, we claim that the envelope
\[
\cP \phi=\sup{}^*\{\psi\in C^0\cap \PSH(L|_M),\,\psi\leq\phi\text{ on }L|_{\partial M}\}
\]
is our desired relatively maximal, continuous metric on $L|_M$ which coincides with~$\phi$ on $L|_{\partial M}$. Because there exists a continuous subsolution, it follows by definition that~$\cP \phi$ is relatively maximal; furthermore, because a subsolution is a candidate $\psi$ to the envelope which coincides with $\phi$ on $L|_{\partial M}$, it also follows that $\cP\phi$ has the correct boundary values. We are therefore left to show continuity of $\cP\phi$.

We begin with a continuity estimate near the boundary. Having fixed a reference smooth, strictly psh metric $\refmetric$ on $L$ (which we assume to exist, up to possibly restricting to a smaller relatively compact open subset inside $\Omega$), and setting $\omega=dd^c\refmetric$, we can see any candidate $\psi$ for the envelope $\cP\phi$ as a continuous $\omega$-psh function $g=\psi-\refmetric$. Fix such a $g$, and set
\[
f_0=\phi-(\refmetric)|_{\partial M}.
\]
We now argue as in \cite[\S 4.2(a)]{phongsturmeq}. Since $dd^c g\geq -\omega$, the Laplacian $\Delta_\omega g$ of $g$ with respect to $\omega$ is bounded below by $-d-1$. Let $f$ be the (continuous) solution on $\bar M$ to the Dirichlet problem
\[
\Delta_\omega f +(d+1) = 0,\,f|_{\partial M}=f_0.
\]
We then have that $\Delta_\omega(g-f)\geq 0$, which implies by the maximum principle that
\[
\sup_{M} \,(g-f) = \sup_{\partial M}\,(g-f),
\]
while this supremum is nonpositive since $\psi=g+\refmetric$ is a candidate for the envelope~$\cP\phi$. We then have that $g\leq f$ on all of $\bar M$, and this is true for any candidate $\psi$, so that $\cP\phi\leq \refmetric + f$ on $\bar M$.

We now look at the continuity on $M$. Let $\tilde\phi$ denote a continuous psh extension of $\phi$ to $L|_M$. We fix $\varepsilon>0$, and define $U=\{\cP\phi<\tilde\phi+\varepsilon\}$. This is the complement of a compact set $C_\varepsilon$, with $\partial\Omega\subset C_\varepsilon$, and $C_\varepsilon$ shrinks as $\varepsilon\to 0$. By regularization (\eg \cite[Th.\,3.8]{boul2}), we can find a sequence $\psi_k\in C^0\cap \PSH(L|_{U})$ which decreases to $\cP\phi$. Now, by Dini, using compactness, we have that $U$ is covered by finitely many of the $U_k=\{\psi_k<\tilde\phi+\varepsilon\}$ (since such inequality holds, for all $z\in M$, and for all large enough $k_z$). In particular, for large enough $k$, one has $\psi_k<\tilde\phi+\varepsilon$. We now define $\tilde\psi_k:=\max(\psi_k-\varepsilon,\tilde\phi)$, which is defined on all of $M$. For all $k$, $\tilde\psi_k$ is continuous, as $\psi_k$ is continuous away from the boundary and $\tilde\phi$ is continuous everywhere (in particular, near and up to the boundary). Furthermore, $\tilde\psi_k$ is equal to $\phi$ on the boundary, so~that
\[
\psi_k-\varepsilon\leq\tilde\psi_k\leq \cP\phi \leq \cP\phi^* \leq \psi_k.
\]
This implies that $\psi_k$ converges uniformly to $\cP\phi$, \ie $\cP\phi$ is continuous on $M$. Furthermore, since:
\begin{enumerate}
\item $\cP\phi\leq \refmetric + f$, as at the end of the first point of the proof;
\item $\refmetric + f$ converges continuously to $\phi$ near the boundary, and $\cP\phi$ is continuous on $M$;
\item there exists a psh extension $\tilde\phi$, ensuring that $\cP\phi= \phi$ on the boundary,
\end{enumerate}
then $\cP\phi$ is continuous up to the boundary.

\subsubsection*{Proof of the second point: construction of a subsolution under the hyperconvexity assumption}
The second point of the theorem will follow from a more general principle: consider the class $\cC(L|_{\partial M})$ consisting of continuous, fiberwise psh metrics on $L|_{\partial M}$ admitting a continuous psh extension to all of $L|_M$. We claim that this class is stable under uniform limits. Indeed, pick such a metric $\psi$, and consider the envelope
\[
\cP \psi=\sup{}^*\{\psi'\in C^0\cap \PSH(L|_M),\,\psi'\leq\psi\text{ on }L|_{\partial M}\},
\]
the star denoting usc regularization. One sees that $\cP$ is increasing (\ie $\psi'\leq \psi$ implies $\cP\psi'\leq\cP\psi$ and that, given a constant $C\in\bbr$, $\cP(\phi+C)=\cP(\phi)+C$, from which it follows formally that the mapping $\cC(L|_{\partial M})\ni \phi\mto \cP\phi$ is continuous under uniform convergence.

To prove the second point, we therefore have to show that there exists a sequence $\phi_k\in\cC(L|_{\partial M})$ converging uniformly to our boundary data $\phi$. We proceed by Bergman kernel approximation, \ie we will construct this sequence using the data of relative sections of $L$. Since $L$ is $\pi$-ample, the sheaves $\pi_*(kL)$ are locally free for all $k$ large enough, and correspond to the sections of a vector bundle $E_k$ whose fibers are the $H^0(kL_z)$, $z\in M$. We then have a continuous collection of Hermitian metrics $h_k=(h_{k,z})_z$ on $E_k|_{\partial \Omega}$, which (identified with their associated Hermitian norms $\norm_{h_{k,z}}$) are given by
\[
\|s_z\|_{h_{k,z}}^2=\int_{X_z}|s_z|^2 e^{-k\phi_z}\,d\omega_z,
\]
where $\omega_z$ is the restriction to $X_z$ of a Kähler form on $X$, and $s=(s_z)_z$ is a section of $E_k$. We pick a sequence of smooth families of Hermitian metrics $(h_{k,j})_j$ on $\pi_*(kL)$ so that $h_{k,j}\to h_k$ uniformly on $\pi_*(kL)|_{\partial \Omega}$. We recall that, to a Hermitian norm $h$ on the space of sections of $kL_z$, we can associate a smooth, strictly plurisubharmonic metric on $L_z$ given by
\[
FS(h)=k\mi\log\sum_i |s_i|^2,
\]
where $(s_i)_i$ is an orthonormal basis for $h$, and the metric so obtained is independent of the choice of such a basis. Using this construction, we obtain collections of smoothly varying metrics
\[
\phi_{k,z}=FS(h_{k,z}),\,\phi_{k,j,z}=FS(h_{k,j,z})
\]
on $L$. Since they are fiberwise \textit{smooth} and \textit{strictly psh} (both of which are necessary conditions for the following argument), we may compensate for the lack of plurisubharmonicity in the direction of $z$, by pulling back a high enough multiple $m_{k,j}\pi^*\rho$ of the defining function $\rho$ of $\Omega$, which as we recall vanishes on the boundary of $\Omega$. We~therefore have a continuous psh extension
\[
\phi_{k,j}+m_{k,j}\pi^*\rho
\]
of $(\phi_{k,j})|_{\partial M}$. This implies that $\phi_{k,j}\in \cC(L|_{\partial M})$; furthermore, $(\phi_{k,j})|_{\partial M}\to \phi_k$ uniformly which implies that $\phi_k\in\cC(L|_{\partial M})$. Now, by Bergman kernel asymptotics in families, as in \eg \cite[Th.\,4.1.1]{mamarinescu}, the $\phi_k$ themselves converge uniformly \textit{and increasingly} to $\phi$, which implies $\phi\in\cC(L|_{\partial M})$, concluding the proof.
\end{proof}
\begin{lemma}\label{lem_relmaxconditions}
In the setting of Theorem \ref{thm_existencemaximal}, let $\phi$ be a continuous psh metric on~$L$. Then, the following are equivalent:
\begin{enumerate}
\item $\phi$ is relatively maximal;
\item $(dd^c\phi)^{d+1}=0$, where $d$ is the dimension of the fibers of $\pi$;
\item over each relative open subset $\pi\mi(U)$, $\phi$ equals its Perron-Bremmermann envelope
\[
\cP_U \phi=\sup{}^*\{\psi\in C^0\cap \PSH(L|_{\pi\mi(U)}),\,\psi\leq\phi\text{ on }L|_{\pi\mi(\partial \bar U)}\}.
\]
\end{enumerate}
\end{lemma}
\begin{proof}
We start with (1) $\Leftrightarrow$ (2). Recall that $\phi$ being relatively maximal means that, for all relative open subsets $\pi\mi(U)$, $\phi|_{\pi\mi(U)}$ is maximal in the sense of Sadullaev. Because maximality in the sense of Sadullaev is a local property, it is enough to check it on any open covering of $\pi\mi(U)$, so that we may work locally and apply \cite[Th.\,2.20]{complexmaeq}. Thus we find that $(dd^c\phi)^{d+1}=0$ on $\pi\mi(U)$ if and only if $\phi|_{\pi\mi(U)}$ is maximal in the sense of Sadullaev, \ie if and only if $\phi$ is relatively maximal.

(1) $\Leftrightarrow$ (3) follows from the definition of (relative) maximality and of the envelope~$\cP_U$; one simply has to note that there always exists a candidate for the envelope~$\cP_U(\phi)$, given by the restriction of $\phi$ itself to $\pi\mi(U)$.
\end{proof}

We now characterize relatively maximal metrics of relative finite energy.
\begin{prop}\label{prop_harmonicitymaximaldiscs}
Let $\pi:X\to Y$ be a proper holomorphic submersion. Let $\phi$ be a metric in $\cE^1_{X/Y}(L)$. Let $U$ be a relatively compact, smooth open subset of $Y$, such that $\pi$ is a submersion above $\partial U$. Then, $\phi$ is relatively maximal on $U$ if and only if the restriction of $\langle \phi^{d+1}\rangle_{X/Y}$ to $U$ has zero curvature.
\end{prop}
\begin{proof}
Assume $\phi$ to be relatively maximal on $U$. If $\phi$ is continuous, it follows from Lemma \ref{lem_relmaxconditions} that $(dd^c\phi)^{d+1} = 0$ on $U$, so that by \eqref{eq_ddc}, it follows that $dd^c \langle \phi^{d+1}\rangle\equiv0$. The non-continuous case follows from regularization on $U$: pick a sequence of continuous metrics $\phi_k$ decreasing to $\phi$ on $U$; by Theorem \ref{thm_existencemaximal}, there exists a continuous, relatively maximal psh metric $\Phi_k$ coinciding with $\phi_k$ on $L|_{\pi\mi(\partial U)}$. By maximality, the sequence $\Phi_k$ necessarily converges to $\phi$ (since $\phi$ is assumed to be relatively maximal), and continuity of Deligne pairings along decreasing nets (which is implicit from the proof of Theorem \ref{thm_eonepairing}) ensures $\langle \phi^{d+1}\rangle_{X/U}$ to have zero curvature.

Conversely, assume $\langle \phi^{d+1}\rangle_{X/U}$ to have zero curvature. In the continuous case, \hbox{using}~\eqref{eq_ddc} again, it follows that $(dd^c\phi)^{d+1}=0$, as it is a positive current. In the general case, we again proceed base-locally, and approximate $\phi$ on the preimage of a~relatively compact open subset $U$ via a decreasing sequence of continuous psh metrics $k\mto \phi_k$. Let $\Phi_k$ be for each $k$ the unique continuous and relatively maximal metric on $U$ with prescribed boundary condition $\phi_k|_{\pi\mi(\partial U)}$, given by Theorem \ref{thm_existencemaximal}. Let $\Phi$ denote the limit of the decreasing sequence $k\mto \Phi_k$, which is relatively maximal. By continuity of the Deligne pairing along decreasing nets, this sequence also defines a decreasing sequence of zero curvature metrics $\langle \Phi_k^{d+1}\rangle_{X/U}$
which has to converge to the metric $\langle \Phi^{d+1}\rangle_{X/U},$ which is a zero curvature metric $\tilde\phi$ on $U$, coinciding on $\partial U$ with $\langle \phi^{d+1}\rangle_{X/U}$. Since $\langle \phi^{d+1}\rangle_{X/U}$ also has zero curvature, we must have
\[
\langle \Phi^{d+1}\rangle_{X/U}=\langle \phi^{d+1}\rangle_{X/U}
\]
on all of $U$. Fix $z$ in $U$, and note that this implies
\[
E(\Phi_z)=E(\phi_z),
\]
while by relative maximality of $\Phi$, $\phi_z\leq \Phi_z$, which implies $\Phi_z=\phi_z$, thus concluding our proof.
\end{proof}

\section{Finite-energy metrics over degenerations}\label{sect_sect15}

\subsection{Analytic models and degenerations}

We now turn to our main setting. We~will consider the base $Y$ to be the punctured unit disc, and we will assume that our family degenerates (meromorphically) as one approaches zero.

\begin{defi}
Consider a holomorphic submersion $\pi:X\to\pdisc$ with projective fibers, and a relatively ample line bundle $L$ on $X$. An analytic model (or~simply a~model) of $X$ is the data of:
\begin{enumerate}
\item a normal complex analytic space $\cX$, together with a flat, proper holomorphic morphism $\pi:\cX\to\overline\bbd$;
\item an isomorphism $X\simeq \pi\mi(\overline\bbd^*)$.
\end{enumerate}
An analytic model of $(X,L)$ is the data of:
\begin{enumerate}
\item an analytic model $\cX$ of $X$;
\item an ample line bundle $\cL$ over $\cX$;
\item an isomorphism between the restriction of $\cL$ to $\pi\mi(\overline\bbd^*)$, and $L$.
\end{enumerate}
We define a degeneration (or a degeneration with meromorphic singularities) to be a morphism $\pi:X\to\pdisc$ as above, such that there exists an analytic model of $(X,L)$.
\end{defi}

\begin{example}
This construction specializes to the following well-known cases:
\begin{itemize}
\item if all the fibers of $X$ are isomorphic to $M$, a model $\cX$ can simply be viewed as a compactification of an isotrivial degeneration of $M$;
\item if the above condition holds, and furthermore the isomorphism is generated by a $\bbc^*$-action, this is simply a (real) one-parameter degeneration of $(M,L|_M)$, \ie a test configuration for $(M,L|_M)$.
\end{itemize}
\end{example}

The central fiber of a model of $X$ is the space $\cX_0=\pi\mi(\{0\})$. If the degeneration $X\to\overline\bbd^*$ is isotrivial, we say that $M$, the fiber over $1$, is the generic fiber of $X$.

\subsection{Generalized slopes and Lelong numbers}\label{sect_lelong}

As we will be working with (generalized) subharmonic functions on the base $\pdisc$, we will often have to work with some notions of Lelong numbers. We review some (old and new) facts in this section.

\begin{defi}We say that a subharmonic function $f$ on $\bbd^*$ has logarithmic growth (near zero) if there exists a real number $a$ such that $f(z)+a\log|z|$ is bounded above near zero.
\end{defi}

Much as one can define the slope at infinity of a convex function, we can define the (possibly signed) Lelong number of a function with logarithmic growth, as follows.
\begin{defi}Given a subharmonic function $f$ with logarithmic growth on $\pdisc$, we define its generalized slope (or generalized Lelong number at zero) to be the value
\[
\nu_0(f):=\lim_{r\to 0}\frac{\sup_{|z|=r} f(z)+a\log|z|}{\log r}-a,
\]
where $a$ is a real number such that $f+a\log|z|$ is bounded near zero. In particular, $\nu_0(f)$ is independent of the choice of such an $a$.
\end{defi}
The interested reader may find more details on the properties of Lelong numbers in \cite[\S 2]{boul2}.

\begin{example}\label{ex_minuslelong}If $f$ is a $\bbs^1$-invariant subharmonic function, we can identify it with a convex function $\tilde f$ of $t=\log r\in(-\infty,0]$, in which case the Lelong number of $f$ simply computes the slope of $\tilde f$ at infinity:
\[
\nu_0(f)=\lim_{t\to -\infty}\frac{\tilde f(t)}{t}.
\]
In our case, we will be interested in generalizing slopes at infinity of non-decreasing convex functions $\tilde g:[0,\infty)\to\bbr$, identified with subharmonic functions on the disc via $g(z)=\tilde g(-\log|z|)$. A computation shows that the slope at infinity is then captured by \textit{minus} the Lelong number of $g$:
\[
\lim_{u\to\infty} \frac{\tilde g(u)}{u}=-\nu_0(g).
\]
\end{example}

\begin{remark}As a consequence of Harnack's inequality (\cite[Cor.\,1.9]{boul2}) $\nu_0(f)$ may equivalently be computed using the integrals
\[
\fint_{|z|=r}f(z)\,dz=(2\pi r)\mi\int_{|z|=r}f(z)\,dz
\]
in place of the suprema.
\end{remark}

We conclude this section with an estimate that will be useful later on. It is the subharmonic version of the estimate
\[
f(s)\leq \lim_{t\to -\infty}\frac{f(t)}{t} s + f(0)
\]
for a convex function $f$ on $(-\infty,0]$.
\begin{lemma}\label{lem_lelongestimate}
Let $f$ be a subharmonic function with logarithmic growth on $\pdisc$. Then, for all $z$, we have
\[
f(z)\leq \nu_0(f)\cdot \log|z|+\sup_{|z|=1} f(z).
\]
\end{lemma}
\begin{proof}
Set $g(s):=\sup_{|z|=e^s} (f(z)+a\log|z|)$. Let $z\in\bbd^*$ and $0<r<|z|$, then set $t=\log r$. Since $g$ is convex on $(-\infty,0]$ (in particular on $[t,0]$), we have that
\begin{align*}f(z)+a\log|z|\leq g(\log |z|)&\leq \Bigl(1-\frac{\log|z|-t}{-t}\Bigr)g(t)+\frac{\log|z|-t}{-t}g(0)\\
&\leq \frac{\log|z|g(t)}{t}+\frac{\log|z|g(0)}{-t}+g(0).
\end{align*}
Taking the limit as $t=\log r\to-\infty$ we then find
\[
f(z)+a\log|z|\leq \nu_0(f)\cdot\log|z|+a\log|z|+g(0),
\]
which is exactly the estimate in the statement of the lemma.
\end{proof}

\subsection{Plurisubharmonic metrics on degenerations}

Fix now a degeneration $\pi:X\to \overline\bbd^*$, endowed with a relatively ample line bundle $L$. We will take our interest to plurisubharmonic metrics on $L$, and in particular their singularities. However, a general psh metric on a degeneration can behave very poorly near the singularity, even though we have assumed existence of an analytic model of $X$. Thus, we need to enforce a rather natural growth condition on such psh metrics, akin to that of linear growth for geodesic rays.

\begin{defi}\label{def_tame}
We say that a psh metric $\phi$ on $L$ has \textit{logarithmic growth} if there exists a model $(\cX,\cL)$ of $(X,L)$ such that $\phi$ extends as a psh metric on $\cL$.
\end{defi}

\begin{defi}\label{def_psh2}We will write $\PSH(L)$ for the space of psh metrics of logarithmic growth on $L$, such which do not restrict to $-\infty$ on a given fiber. If it comes to be necessary, we will rather write $\PSH(X,L)$ when considering the space of (non-necessarily of logarithmic growth) psh metrics on $L$.
\end{defi}

We will soon show that $\PSH(L)$ has many desirable properties. We will also shortly explain our terminology. We begin with the following result:
\begin{lemma}\label{lem_extpshtame}Given a psh metric $\phi$ on $L$, the following are equivalent:
\begin{enumeratei}
\item\label{lem_extpshtamei} $\phi$ has logarithmic growth, \ie there exists a model $(\cX,\cL)$ such that $\phi$ extends to a psh metric on $\cL$;
\item\label{lem_extpshtameii} for all models $(\cX,\cL)$ of $(X,L)$, there exists a constant $c=c(\cX,\cL)$ such that $\phi+c\cdot \log|z|$ extends to a psh metric on $\cL$;
\item\label{lem_extpshtameiii} there exists a model $(\cX,\cL)$ and a smooth metric $\refmetric$ on $\cL$ such that
\[
\rho^*\phi(z) \leq \refmetric(z) + O(\log |z|)
\]
as $z\to 0$, where $\rho$ denotes the isomorphism between $X$ and $\cX-\cX_0$;
\item\label{lem_extpshtameiv} for all models $(\cX,\cL)$ of $(X,L)$ and all smooth metrics $\refmetric$ on $\cL
$, \eqref{lem_extpshtameiii} holds.
\end{enumeratei}
\end{lemma}
\begin{proof}
By classical results of pluripotential theory, $\eqref{lem_extpshtamei}\Leftrightarrow\eqref{lem_extpshtameiii}$ and $\eqref{lem_extpshtameii}\Leftrightarrow\eqref{lem_extpshtameiv}$. Since $\eqref{lem_extpshtameiv}\Rightarrow\eqref{lem_extpshtameiii}$ is immediate, we only need to prove $\eqref{lem_extpshtameiii}\Rightarrow\eqref{lem_extpshtameiv}$. Assume that
\[
\rho^*\phi(z) \leq \phi^\mathrm{ref}_\cL(z) + O(\log |z|)
\]
for a smooth reference metric $\phi^\mathrm{ref}_\cL$ on $\cL$. Pick another model $(\cY,\cM)$ together with a smooth metric $\phi^\mathrm{ref}_\cM$. Note that the equation above holds if and only if the same equation holds for the pullbacks of $\phi^\mathrm{ref}_\cL$ and $\rho^*\phi$ to a higher model. Thus, we pick a model $(\cZ,\cN)$ dominating both via $\pi_\cX:\cZ\to \cX$, $\pi_\cY:\cZ\to\cY$. There exists a unique Cartier divisor $D$ supported on the special fiber $\cZ_0$ such that\vspace*{-3pt}
\[
\pi_\cX^*\cL+D=\pi_\cY^*\cM,
\]
and given a local equation $f_D$ for $D$, we have\vspace*{-3pt}
\[
\pi_\cX^*\phi_\cL \leq \pi_\cY^*\phi^\mathrm{ref}_\cM - \log|f_D| + O(1)\leq \pi_\cY^*\phi^\mathrm{ref}_\cM + O(\log|z|).
\]
Thus,\vspace*{-3pt}
\[
\pi_\cX^*\rho^*\phi\leq \pi_\cX^*\phi^\mathrm{ref}_\cL + O(\log|z|)\leq \pi_\cY^*\phi^\mathrm{ref}_\cM + O(\log|z|),
\]
as desired.
\end{proof}

\begin{remark}The above result shows that one could equivalently define our growth condition using some fixed reference data $(\cX_\mathrm{ref},\cL_{\mathrm{ref}})$, using \eg point \eqref{lem_extpshtameii}. In the isotrivial case, there furthermore exists some very natural reference data: the ``trivial model'' given by the product family of the generic fiber with the whole disc.
\end{remark}

\begin{example}
Let $[0,\infty)\ni t\mto \phi_t$ be a ray of psh metrics on an ample line bundle $L$ over a fixed variety $X$. It may be identified as a psh metric $\Phi$ over the trivial model $(X\times\overline\bbd^*,L\times\overline\bbd^*)$, by setting $\Phi_z=\phi_{-\log|z|}$. In this case, the logarithmic growth condition is merely the usual linear growth condition on psh rays.
\end{example}

We then have as an immediate corollary:
\begin{corollary}The space $\PSH(L)$ is stable under limits of decreasing nets, finite maxima, and addition of constants. It is furthermore the smallest such set containing all psh metrics on $L$ which admit a locally bounded extension to some model $(\cX,\cL)$ of $(X,L)$.
\end{corollary}
\begin{proof}
All of those properties are seen to preserve characterization \eqref{lem_extpshtameiv} above, having fixed some reference model. To show that it is the smallest set closed under those operations, only the statement about decreasing nets could a priori be delicate. Given a metric $\phi\in\PSH(L)$, \eqref{lem_extpshtamei} shows that it extends as a genuine metric on some model $(\cX,\cL)$, and Demailly's regularization theorem yields a decreasing sequence of smooth (in particular locally bounded) psh metrics decreasing to the extension of $\phi$, which shows in particular that $\phi$ belongs to the closure of the set of locally bounded psh metrics on $\cL$, proving our result.
\end{proof}

\subsection{The main setting, and some important examples}

We begin with some notation. Let $\pi:X\to\pdisc$ be a degeneration together with a relatively ample line bundle~$L$. We now, and for the remainder of this article, fix some reference boundary data $\phi_\partial$, which is the restriction to the boundary $\pi\mi(\bbs^1)$ of a smooth psh metric on $L$. This is a minor distinction which will allow us to later obtain a genuine metric structure on a particular subspace of $\PSH(L)$, rather than a pseudometric structure, and therefore we will define
\[
\cE^1(L)=\cE^1_{X/\pdisc}(L)\cap\{\phi\in\PSH(L),\,\phi\text{ has boundary data }\phi_\partial\}
\] to be the space of fiberwise finite-energy metrics in $\PSH(L)$ (with the correct boundary data). We also set
\[
\hat\cE^1(L)=\{\phi\in\cE^1(L),\,\phi\text{ is relatively maximal}\}.
\]

\begin{example}Although those are seemingly restrictive conditions, they are in fact general enough to encompass the study of maximal geodesic rays. Let $(X,L)$ be a product family $(M\times\pdisc,L_M\times\pdisc)$. For a $\bbs^1$-invariant metric $\phi$ on $L_M\times\pdisc$, seen as a ray $[0,\infty)\ni t\mto \phi_t,$
\begin{enumerate}
\item being in $\PSH(L)$ corresponds to the usual linear growth condition;
\item being relatively maximal corresponds to being a geodesic ray in the sense of \cite{bbj};
\item being in $\cE^1(L)$ corresponds to having fiberwise finite-energy and linear growth, as well as having a fixed value $\phi_\delta\in\PSH(M)$ at $t=0$;
\item therefore, belonging to $\hat\cE^1(L)$ corresponds to being a fiberwise finite-energy geodesic rays with linear growth emanating from a given metric $\phi_\delta$---exactly the space of rays $\cR^1(L)$ considered in \cite{darlurays}.
\end{enumerate}
\end{example}

\begin{example}[Relative dimension zero, part 1]\label{exa_reldim}Consider the case of relative dimension zero with a (trivial) line bundle $L$ over $X\simeq \bbd^*$. Then,
\begin{enumerate}
\item $\PSH(L)$ corresponds to the set of subharmonic functions with logarithmic growth on $\bbd^*$;
\item the class of relatively maximal metrics in $\PSH(L)$ corresponds to the class of harmonic functions;
\item $\cE^1(L)$ corresponds to finite-valued subharmonic functions on the punctured disc with logarithmic growth (and having determined boundary values on $\bbs^1$);
\item finally, $\hat\cE^1(L)$ corresponds to finite-valued harmonic functions on the punctured disc with logarithmic growth and with the fixed boundary values on $\bbs^1$.
\end{enumerate}

It is well-known that any harmonic function on the punctured disc decomposes as a sum of a multiple of $\log|z|$ and the real part of an analytic function. This is where our general setting starts diverging from the better-behaved $\bbs^1$-invariant. Indeed, by \cite[Prop.\,4.1]{bbj}, for rays of metrics of finite energy, maximality implies linear growth. However, in our case, maximality plus finite energy no longer implies logarithmic growth, since there exist harmonic functions on the punctured disc that do not have logarithmic growth at zero (\eg the real part of $z\mto e^{1/z}$).

Assuming logarithmic growth, we then have a full description of $\hat\cE^1(L)$ in relative dimension zero, since we then see that any (finite-valued) harmonic function with logarithmic growth has to be of the form $c\cdot \log|z|+H(z)$, where $H(z)$ is the solution of the generalized Dirichlet problem over the whole disc with the given boundary data. In particular, it is an affine space isomorphic to $\bbr$. This agrees with the radial case, where $\cE^1(L)$ is simply the set of affine functions on $[0,\infty)$ emanating from the same point, which is isomorphic to the set of possible slopes.
\end{example}

\begin{example}[Relative dimension zero, part 2]\label{exa_reldim2}
Note that, in this setting, the existence of a model for $(\pdisc,L)$ means that we can pick a trivialization $\tau$ of $\cL$ over $\disc$, which allows us to identify a metric $\phi\in\PSH(L)$ (extended to $\cL$ via the logarithmic growth condition) with the function
\[
u=-\log|\tau|_\phi
\]
on $\disc$. By the discussion above, if $dd^c\phi=0$, then $u$ decomposes as
\[
u(z)=c\cdot \log|z|+H(z),
\]
where $H$ is bounded on $\disc$. This decomposition (in particular, $c$ and $H$) depends on $\tau$; but the fact that $\phi$ can be decomposed in any trivialization in such a way does not.
\end{example}

This is a nice model case for us, because the Deligne pairing construction (in our setting of fibrations over $\pdisc$) naturally gives line bundles over $\pdisc$, as we see in action now.

\begin{corollary}\label{coro_trivialization}The relative maximality condition for metrics in $\cE^1(L)$ can be pushed forward to the base via the Deligne pairing, \ie we have a well-defined map
\[
\hat\cE^1(L)\to \hat\cE^1(\langle L^{d+1}\rangle_{X/\pdisc}).
\]
Furthermore, a metric $\phi\in\cE^1(L)$ belongs to $\hat\cE^1(L)$ if and only if, for any model $(\cX,\cL)$ of $(X,L)$ and any trivialization of the Deligne pairing $\langle \cL^{d+1}\rangle_{\cX/\disc}$, denoting $u=-\log|\tau|_{\langle \phi^{d+1}\rangle_{\cX/\disc}}$, one has
\[
u(z)= c\cdot \log|z|+H(z),
\]
where $c$ is a real constant and $\mathrm{H}$ is a harmonic function on $\disc$ depending only on $\tau$ and the boundary data.
\end{corollary}
\begin{proof}
The map above is naturally given by $\phi\mto \langle \phi^{d+1}\rangle_{X/\pdisc},$ in which case both statements are corollaries of Proposition \ref{prop_harmonicitymaximaldiscs} and the two examples above.
\end{proof}

\subsection{Metrization}

By the work of Darvas-Lu (\cite{darlurays}), it is known that one can endow the space of maximal psh rays with a metric structure, given by
\[
\hat d_1(\phi_0,\phi_1)=\lim_{t\to\infty}\frac{d_1(\phi_{0,t},\phi_{1,t})}{t}.
\]
As consequence of our previous results, we will define in this section a metric structure on our space $\hat\cE^1(L)$ of ``generalized rays'', which will be given by
\[
\hat d_1(\phi_0,\phi_1):=-\nu_0(d_1(\phi_0,\phi_1)),
\]
where $d_1(\phi_0,\phi_1)$ is the function $z\mto d_1(\phi_{0,z},\phi_{1,z})$. The reason for the minus sign appearing is explained in Example \ref{ex_minuslelong}. In the following sections, we will show that this structure furthermore satisfies some good properties, namely completeness and geodesicity.
\begin{theorem}\label{thm_1mainthmsepdiscs}
The space $\hat\cE^1(L)$ can be endowed with a metric space structure, defined by the generalized slope $\hat d_1(\phi_0,\phi_1)=\nu_0(d_1(\phi_0,\phi_1))$ for any $\phi_0$, $\phi_1\in\hat\cE^1(L)$.
\end{theorem}

Naturally, this suggests that the $d_1$-distance is subharmonic with logarithmic growth along metrics in $\hat\cE^1(L)$, a fact that we prove now.
\begin{prop}\label{prop_distancesubharmonic}
Let $\phi_0$, $\phi_1\in\hat\cE^1(L)$. Then, the map
\[
z\mto d_1(\phi_{0,z},\phi_{1,z})
\]
is subharmonic with logarithmic growth on $\bbd^*$.
\end{prop}
\begin{proof}
By the formula for $d_1$,
\[
d_1(\phi_{0,z},\phi_{1,z})=\langle\phi_{0,z}^{d+1}\rangle+\langle\phi_{1,z}^{d+1}\rangle-2\langle P(\phi_{0,z},\phi_{1,z})^{d+1}\rangle.
\]
By Proposition \ref{prop_harmonicitymaximaldiscs}, the first two metrics on the right-hand side have zero curvature, therefore we are left to show that the metric $\langle P(\phi_{0},\phi_{1})^{d+1}\rangle_{X/\pdisc}$ is superharmonic. We pick any zero curvature metric $\refmetric$ on $\bbd^*$, and note that $\langle P(\phi_{0},\phi_{1})^{d+1}\rangle_{X/\pdisc}$ is superharmonic if and only if $\langle P(\phi_{0},\phi_{1})^{d+1}\rangle_{X/\pdisc}-\refmetric$ is a superharmonic function. Fix $a\in\bbd^*$ and let $r>0$ be such that $\bbd(a,r)=\{|z-a|\leq r\}\subset \bbd^*$. Let $\psi$ be the relatively maximal psh metric on $\bbd(a,r)$ and with boundary data
\begin{equation}\label{eq_lem_dist_1}
\psi(z)=P(\phi_{0,z},\phi_{1,z}),\,\forall z\in S(a,r).
\end{equation}
Such a metric is given by Theorem \ref{thm_existencemaximal}. We now deduce the two following facts:
\begin{enumeratei}
\item\label{enum:i}
by maximality of $\psi$, if follows from Proposition \ref{prop_harmonicitymaximaldiscs} that $z\mto \langle \psi^{d+1}\rangle_{X/\pdisc}$ has zero curvature;
\item\label{enum:ii}
since on the boundary $S(a,r)$ we have $\psi(z)\leq \phi_{0,z},\phi_{1,z}$, and $\phi_{0}$, $\phi_{1}$ are relatively maximal, we have
\[
\psi_z\leq \phi_{0,z},\phi_{1,z}
\]
for all $z\in\bbd(a,r)$, thus $\psi_z\leq P(\phi_{0,z},\phi_{1,z})$ and finally
\[
\langle\psi_z^{d+1}\rangle\leq \langle P(\phi_{0,z},\phi_{1,z})^{d+1}\rangle
\]
by monotonicity of the Monge-Ampère energy.
\end{enumeratei}
Using (\ref{eq_lem_dist_1}), \eqref{enum:i}, and \eqref{enum:ii} in that order, we find:
\begin{align*}
\fint_{S(a,r)}\langle P(\phi_{0,z},\phi_{1,z})^{d+1}\rangle - \phi_{\mathrm{ref},z}&=\fint_{S(a,r)} \langle \psi_z^{d+1}\rangle - \phi_{\mathrm{ref},z}\\
&=\langle \psi_a^{d+1}\rangle - \phi_{\mathrm{ref},a}\\
&\leq \langle P(\phi_{0,a},\phi_{1,a})^{d+1}\rangle - \phi_{\mathrm{ref},a}.
\end{align*}
As the inequality is true for all $a$, our metric $\langle P(\phi_{0},\phi_{1})^{d+1}\rangle_{X/\pdisc}$ is then superharmonic (because superharmonic functions are characterized by the ``reverse'' mean-value inequality, \ie the value of a superharmonic function at a point is no smaller than the mean of its values on a circle or a ball centered at the given point). We now show that there exists a real number $a\in\bbr$ such that
\[
z\mto d_1(\phi_{0,z},\phi_{1,z})+a\log|z|
\]
is bounded above. By Lemma \ref{lem_extpshtame}(iv), for any model $(\cX,\cL)$ of $(X,L)$, fixing a reference metric $\refmetric\in\hat\cE^1(L)$ which is locally bounded on $\cL$, one has (up to adding large enough constants)
\[
\phi_0\leq\refmetric + c\cdot\log|z|
\]
for some real constant $c$. In this case,
\[
d_1(\phi_{0,z}-c\cdot \log|z|,\phi_{\mathrm{ref},z})=\langle(\phi_{0,z}-c\cdot \log|z|)^{d+1}\rangle-\langle\phi_{\mathrm{ref},z}^{d+1}\rangle,
\]
and the term on the right-hand side is a harmonic function with logarithmic singularities at the origin, so that subtracting constants the result also holds for $z\mto d_1(\phi_{0,z},\phi_{\mathrm{ref},z})$. Proceeding similarly for $\phi_1$, our result then follows from the triangle inequality.
\end{proof}

Finally, we note an immediate consequence of Lemma \ref{lem_lelongestimate} together with the previous proposition \ref{prop_distancesubharmonic}.
\begin{lemma}\label{lem_unifcontrol}
Let $\phi_0$, $\phi_1\in\hat\cE^1(L)$. Then, for all $z$ on the base, we have
\[
d_1(\phi_{0,z},\phi_{1,z})\leq \hat d_1(\phi_0,\phi_1)(-\log|z|).
\]
In particular, $\hat d_1(\phi_0,\phi_1)$ is nonnegative.
\end{lemma}
\begin{remark}Had we not fixed boundary data, we would have an additional error term in the above expression, corresponding exactly to the supremum of $z\mto d_1(\phi_{0,z},\phi_{1,z})$ for $z\in\bbs^1$.
\end{remark}

We are now equipped to endow the space $\hat\cE^1(L)$ with a metric structure.
\begin{proof}[Proof of Theorem \ref{thm_1mainthmsepdiscs}.]
That $\hat d_1(\phi,\phi)=0$ and $\hat d_1(\phi_0,\phi_1)=\hat d_1(\phi_1,\phi_0)$ are immediate statements, and nonnegativity will follow from the triangle inequality and the former statement. Therefore, we must show that for any other $\phi_2\in\hat\cE^1(L)$, we have
\[
\hat d_1(\phi_0,\phi_1)\leq \hat d_1(\phi_0,\phi_2)+\hat d_1(\phi_2,\phi_1).
\]
Let $a_{01}$ be such that $d_1(\phi_{0,z},\phi_{1,z}) + a_{01}\log|z|$ is bounded above on the punctured disc, and define similarly $a_{02}$, $a_{21}$. We have by the triangle inequality of the fiberwise metric $d_1$
\[
d_1(\phi_{0,z},\phi_{1,z})\leq d_1(\phi_{0,z},\phi_{2,z})+d_1(\phi_{2,z},\phi_{1,z})
\]
for all $z$ in $\bbd^*$, and in particular
\[
d_1(\phi_{0,z},\phi_{1,z})+(a_{02}+a_{21})\log|z|\leq d_1(\phi_{0,z},\phi_{2,z})+d_1(\phi_{2,z},\phi_{1,z})
+(a_{02}+a_{21})\log|z|.
\]
Upon taking (negative) Lelong numbers and adding constants, we find\vspace*{-3pt}
\begin{multline*}
a_{02}+a_{21}+\nu_0(d_1(\phi_{0,z},\phi_{1,z})+(a_{02}+a_{21})\log|z|)\\
\leq a_{02} +\nu_0(d_1(\phi_{0,z},\phi_{2,z})+a_{02}\log|z|)+a_{21} +\nu_0(d_1(\phi_{2,z},\phi_{1,z})+a_{21}\log|z|).
\end{multline*}
Since
\[
z\mto d_1(\phi_{0,z},\phi_{1,z})+(a_{02}+a_{21})\log|z|
\]
is bounded above, the previous equation is by the very definition of $d_1$ equivalent to
\[
\hat d_1(\phi_0,\phi_1)\leq \hat d_1(\phi_0,\phi_2)+\hat d_1(\phi_2,\phi_1),
\]
as desired. Finally, assuming $\hat d_1(\phi_0,\phi_1)=0$, Lemma \ref{lem_unifcontrol} shows that we must have $\phi_0=\phi_1$.
\end{proof}

\subsection{Completeness}

We wish to prove the following:
\begin{theorem}\label{thm_1completeness}
The metric space $(\hat\cE^1(L),\hat d_1)$ is complete.
\end{theorem}
In order to prove this, we discuss possible topologies for $\cE^1(L)$.
\begin{remark}[Topologies on $\cE^1(L)$]\label{rem_topologies}We have already considered the topology of fiberwise $d_1$-convergence on $\cE^1(L)$. There is a yet finer topology, that of locally uniform fiberwise $d_1$-convergence (the $C^0_{X,\mathrm{loc}}$-$d_1$ topology for short), by which $\phi_k$ converges to~$\phi$ if, for all relatively compact open sets $U$ in $X$, $d_1(\phi_{k,z},\phi_z)\to 0$ uniformly in $z$ on $U$. In between the two, there is the topology of ``base-locally'' uniform fiberwise $d_1$-convergence (abbreviated by $C^0_{X/\pdisc,\mathrm{loc}}$-$d_1$), which is the same but over the $\pi\mi(U)$ with $U$ relatively compact open in $\bbd^*$. By Lemma \ref{lem_unifcontrol}, the latter is equivalent to the topology induced by $\hat d_1$ on $\hat\cE^1(L)$.
\end{remark}

\begin{prop}\label{prop_closed}Let $\phi_k$ be a sequence of metrics in $\hat\cE^1(L)$ converging in the $C^0_{X/\pdisc,\mathrm{loc}}$-$d_1$ topology to some metric $\phi\in\cE^1(L)$. Then, $\phi$ belongs to $\hat\cE^1(L)$.
\end{prop}
\begin{proof}
Pick a sequence $k\mto \phi_k\in\hat\cE^1(L)$ and a fixed metric $\phi$ in $\cE^1(L)$. Assume that, for a relatively compact open $U\subset\pdisc$ we have\vspace*{-3pt}
\[
d_1(\phi_{k,z},\phi_z)\to 0
\]
uniformly in $z\in \pi\mi(U)$. Since convergence in Monge-Ampère energy is subordinate to $d_1$-convergence (\cite[Th.\,3.46]{darbook}) we have that\vspace*{-3pt}
\[
\langle\phi_{k,z}^{d+1}\rangle\to \langle\phi_z^{d+1}\rangle
\]
again uniformly in $z$; by maximality, the metrics $\langle\phi_{k}^{d+1}\rangle_{X/\pdisc}$ have zero curvature, and an uniform limit of such has zero curvature again. As having zero curvature is a local property and the $\pi\mi(U)$ cover $X$, we then have that $\langle\phi^{d+1}\rangle_{X/\pdisc}$ has zero curvature on all of $X$. By virtue of being in $\cE^1(L)$, this implies $\phi$ to be relatively maximal by~\ref{prop_harmonicitymaximaldiscs}.
\end{proof}
\begin{proof}[Proof of Theorem \ref{thm_1completeness}.]
Consider a Cauchy sequence $m\mto \phi_m\in\hat\cE^1(L)$. For all $\varepsilon$ and all large enough $m$, $n$,\vspace*{-3pt}
\[
\hat d_1(\phi_m,\phi_n)\leq \varepsilon,
\]
which by Lemma \ref{lem_unifcontrol} implies the individual sequences $m\mto \phi_{m,z}$ to be $d_1$-Cauchy. By completeness of the fiberwise $\cE^1$ spaces (\cite[Th.\,3.36]{darbook}), those sequences $d_1$-converge to a unique finite-energy metric $\phi(z)$, and in fact this convergence is seen to hold base-locally uniformly fiberwise. The mapping
$z\mto \phi(z)$
is therefore a metric in $\hat\cE^1(L)$ by Proposition \ref{prop_closed}.
\end{proof}

\Subsection{Extension of the distance to non-maximal continuous metrics}

In this section, we construct a ``maximal envelope'' map, which will allow us to extend the $d_1$-distance as a pseudodistance to the class of (non-maximal) continuous metrics in $\cE^1(L)$.

\begin{prop}\label{prop_envelope}For all $\phi\in C^0\cap\cE^1(L)$, there exists a unique smallest relatively maximal metric $\hat P(\phi)\in\hat\cE^1(L)$ with $\phi\leq\hat P(\phi)$ and\vspace*{-3pt}
\[
d_1(\phi_z,\hat P(\phi)_z)=o\,(\log|z|)
\]
as $z\to 0$. This defines a natural projection\vspace*{-3pt}
\[
C^0\cap\cE^1(L)\to C^0\cap\hat\cE^1(L).
\]
\end{prop}

Before proving this result, we note this immediate corollary:
\begin{corollary}The mapping
\[
\hat d_1(\phi_0,\phi_1)=\hat d_1(\hat P(\phi_0),\hat P(\phi_1))
\]
defines a pseudodistance on $C^0\cap\cE^1(L)$.
\end{corollary}

\begin{proof}[Proof of Proposition \ref{prop_envelope}.]
Let $\phi\in\cE^1(L)\cap C^0(L)$, and, for all $r\in (0,1)$, let $U_r$ denote the annulus $\{r < |z| < 1\}\subset \pdisc$, and $V_r=\pi\mi(U_r)\subset X$. Let $\phi_r$ be the relatively maximal metric on $V_r$, coinciding with $\phi$ on $\partial V_r$, given by Theorem \ref{thm_existencemaximal}. Fixing $z$ on the base, the sequence $r\mto \phi_{r,z}$ is an increasing sequence of psh metrics in $\cE^1(L_z)$ by relative maximality, just as, for a convex function on $[0,\infty)$ (corresponding to our $\phi$), the sequence of chords $c_t:[0,t]\to \bbr$ (corresponding to our $\phi_r$) joining $f(0)$ and $f(t)$ forms an increasing family $c_t(\lambda)$; we leave the somewhat long but straightforward adaptation of this argument from the convex to the subharmonic case to the interested reader. We claim that the limit family
\[
z\mto \bigl(\lim_{r\to 0}{}^* \phi_{r,z}\bigr)
\]
is the desired envelope $\hat P(\phi)$. Denote this limit $\hat \phi$ for the moment. Fix some $r$. By~construction, $\hat \phi$ restricted to $V_r$ coincides everywhere with its Perron-Bremmermann envelope; furthermore, it is locally bounded (since it is approximable from below). By~the discussion in Section \ref{sect_relmax}, since this holds for all $r$, $\hat \phi$ is relatively maximal. Furthermore, by construction again, it satisfies $\phi\leq\hat\phi$ and is the smallest such relatively maximal metric. We are therefore only left to prove that $d_1(\hat\phi_z,\phi_z)=o(\log|z|)$ as $z\to 0$. As in Corollary \ref{coro_trivialization}, we pick a model $(\cX,\cL)$ of $(X,L)$, and we extend $\langle \phi^{d+1}\rangle$ to the trivializable line bundle $\langle \cL^{d+1} \rangle$ (using the extension of the Deligne pairing from Moriwaki \cite{mor} when the mapping is no longer a submersion). Picking a trivialization $\tau$ allows us to identify the energies $\langle\hat \phi^{d+1}\rangle$ and the $\langle\phi_r^{d+1}\rangle$ with functions $u$ and $u_r$ on $\bbd$ and $U_r$ respectively. By Proposition \ref{prop_harmonicitymaximaldiscs}, those functions are harmonic, and for all $s\in(0,1)$, the functions $u_r$, $r>s$ increase over $\overline{U_s}$ to $u$, which implies the convergence to be uniform (as an increasing sequence of harmonic functions over a compact set). Now, by harmonicity, for $r>s$, the integrals
\[
\fint_{|z|=r}u_s(z)\,dz
\]
are affine functions of $\log r$. Writing
$v=-\log|\tau|_{\langle \phi^{d+1}\rangle},$
we then have
\[
\fint_{|z|=r}u_s(z)\,dz=\frac{\log r}{\log s}\fint_{|z|=s}v(z)\,dz+\Bigl(1-\frac{\log r}{\log s}\Bigr)\cdot \fint_{|z|=1}v(z)\,dz,
\]
(recall how we have defined $\phi_s$ and $u_s$). Taking the limit $s\to 0$ using the uniform convergence discussed above yields
\[
\fint_{|z|=r}u(z)\,dz=-(\log r)\, \nu_0(v) + \fint_{|z|=1}v(z)\,dz.
\]
Taking slopes in this equality, one then finds $\nu_0(v) = \nu_0(u).$
Now, since $\phi\leq\hat\phi$, we~have
\[
d_1(\phi_{z},\hat \phi_z)=u(z)-v(z),
\]
whose slopes we have seen to coincide, proving our statement that $d_1(\hat\phi_z,\phi_z)=o(\log|z|)$. Therefore, $\hat \phi$ is our desired envelope $\hat P(\phi)$.
\end{proof}

\begin{remark}As pointed out by a referee, because of some problems related to convergence of Lelong numbers under decreasing limits, this construction does not a~priori extend to all of $\cE^1(L)$. It can however be extended to slightly larger classes of metrics, such as metrics $\phi$ for which there exists a sequence of metrics $\phi_i\in C^0\cap\cE^1(L)$ such that $\phi_i\to \phi$ pointwise and the $\nu_0(v_i)$ (which one would define as in the previous proof) converge to $\nu_0(v)$. This is for example the case for the class of metrics which can be written as decreasing limits of continuous metrics in $\cE^1(L)$ and such that convergence also holds in the $C^0_{X/\pdisc,\mathrm{loc}}$-$d_1$ topology introduced in Remark \ref{rem_topologies}.
\end{remark}

\section{The non-Archimedean limit}\label{sect_sect2}

We move away from relatively maximal and finite-energy metrics for the moment, and focus on the space $\PSH(L)$. The purpose of this section is to show that there is a natural map from this space to a certain space of non-Archimedean metrics. We~describe the non-Archimedean setting in Sections \ref{subsect_21} and \ref{sect_nappt}. We prove some complex preparations in Section \ref{subsect_23}, then describe the construction in Section \ref{subsect_24}. Finally, in Section \ref{subsect_25}, we show how a certain subclass of metrics behaves under this map.

\subsection{Degenerations as varieties over a discretely valued field}\label{subsect_21}

Any holomorphic degeneration of algebraic objects naturally gives rise to a corresponding object defined over the field $\bbc(\!(t)\!)$ of formal Laurent series. This idea has been used by Berkovich (\cite{berkhodge}), see also \cite{favre}, \cite{bjtrop}, to build so-called \textit{hybrid spaces}, \ie topological objects mixing complex analytic spaces and Berkovich analytic spaces \hbox{defined} over $\bbc(\!(t)\!)$. Although we will not rely on such constructions, Theorem~\ref{th:B} also relates Archimedean and non-Archimedean objects, and thus fits well into this perspective. For clarity, we will from now on write $\K=\bbc(\!(t)\!)$ and $\R=\bbc[\![t]\!]$.

Pick a degeneration $\pi:X\to\pdisc$ and an analytic model $\pi:\cX\to\disc$ of $X$. As $X$ is projective, it can be embedded in some $\bbp^n\times\bbd$, where it is presented by a finite number of homogeneous polynomials with coefficients in the set of holomorphic functions on~$\pdisc$ that are meromorphic at zero. Since this set of functions can be identified with the field $\K$ of complex Laurent series, one can then view $X$ as a variety $X_\K$ over the field $\K$. Similarly, $\cX$ can be presented by finitely many homogeneous polynomials with coefficients in $\cO(\bbd)$, \ie holomorphic functions over the disc, so that it can be identified with a variety $\cX_\R$ over $\R$.

\begin{example}In the case of a trivial degeneration $X\simeq M\times \bbd^*$ for some complex projective manifold $M$, $X$ can be identified with the base change of $M$ to the field~$\K$. In particular, there exists a ``trivial'' algebraic model, defined by taking the base change of $M$ to $\R$, which corresponds to the product analytic family over $\bbd$.
\end{example}

The field $\K$ is a (non-Archimedean) valued field, with valuation
\[
\nu_0(\Sigma a_i t^i)=\min\{i,\,a_i\neq 0\}.
\]
This also defines a valuation on the Noetherian ring $\R$. From the general work of Berkovich (\cite{berko}), one can associate to a scheme $X$ over a valued ring $\R$, in a functorial way, its analytification $X\an$ with respect to the given valuation on the base. The underlying points of this analytification correspond to pairs $(\zeta,\nu_\zeta)$, where~$\zeta$ is a scheme point of $X$, and $\nu_\zeta$ is a valuation on the residue field of $\zeta$, extending the base valuation on $\mathrm{K}$, and the topology is that of pointwise convergence.

In our setting, the Berkovich analytification $X_\K\an$ of $X_\K$ contains an important dense subset: the set of divisorial points $X^{\mathrm{div}}$. It is described as follows. Let $\cX$ be an analytic model of $X$. We can identify the central fiber $\cX_0$ of $\cX$ over $0$ with a Cartier divisor, which we can decompose as a Weil divisor
\[
\cX_0=\sum_i a_i E_i,
\]
with each $E_i$ irreducible. Each component $E_i$ of such a decomposition defines a valuation $\nu_{E_i}$ on $\K(X)$, given as follows: for all $f\in \K(X)$,
\[
\nu_{E_i}(f)=\ord_{E_i}(f)/a_i.
\]
We then define the set of divisorial points of $X_\K\an$ to be the set of all valuations obtained in this manner. In other words, divisorial points are in correspondence with irreducible component in the central fibers of models of $X$.

\subsection{Non-Archimedean plurisubharmonic functions}\label{sect_nappt}

Let $X$ be a degeneration with a line bundle $L$ over $X$. Let $(\cX,\cL)$ be a model of $(X,L)$. In the same way as before, we can see $L$ as a scheme $L_\K$ on $X_\K$, and $\cL$ as a scheme $\cL_\R$ over $\cX_\R$. To $\cL$ one can associate a model metric $\phi_\cL$ on $L_\K\an$, as explained in detail in \cite{bfjsemi}. Such a metric is uniquely characterized as follows: given a Zariski open set $\mathcal{U}\subset\cX_\K$ and a nonvanishing section of the restriction of $\cL$ to $\mathcal{U}$, then we require that $|s|_{\phi_\cL}=1$ on $(\mathcal{U}\cap X_\K)\an$.

\begin{defi}\label{def_defpsh}We say that a model metric $\phi_\cL$ is plurisubharmonic if $\cL$ is relatively nef. Given a metric $\phi$ on $L_\K\an$, we say that it is plurisubharmonic, and we write $\phi\in\PSH(L_\K\an)$ if there exists a sequence of plurisubharmonic model metrics on $L_\K\an$ decreasing to $\phi$.
\end{defi}

Fixing a psh model metric $\phi_\cL$ on $L_\K\an$, one can identify psh metrics on $L_\K\an$ with ``$\cL$-psh'' functions on $X_\K\an$, via $\phi \leftrightarrow \phi-\phi_\cL$. We define more generally the set of $L$-psh functions to be the unions of all $\cL$-psh functions for all nef models $\cL$ of $L$.

We usually endow the space of $\cL$-psh functions with the topology of pointwise convergence on divisorial points, \ie $\phi_k\to\phi$ in $\PSH(\cL_\R\an)$ if and only if, for all $\nu\in X^{\mathrm{div}}$, $\phi_k(\nu)\to\phi(\nu)$. We note (\cite{bfjsemi}) that a non-Archimedean psh function is uniquely determined by its values on $X^{\mathrm{div}}$.

Any vertical ideal sheaf $\mathfrak{a}$ on a model $\cX$ of $X$ defines a function $\log|\mathfrak{a}|$ on $X\an$, via
\[
\log|\mathfrak{a}|(x)=\max \{\log|f(x)|\},
\]
where the $f$ run over a set of local generators for $\mathfrak{a}$. (In particular, any vertical Cartier divisor $D$ on a model defines such a function.) We then have the following crucial result:
\begin{lemma}[{\cite{bfjsemi}}]
Let $(\cX,\cL)$ be a model of $(X,L)$. Let $\mathfrak{a}$ be a vertical ideal sheaf on $\cX$, such that $\cL\otimes\mathfrak{a}$ is globally generated. Then, $\phi_\cL + \log|\mathfrak{a}|$ is a psh metric on $L_\K\an$.
\end{lemma}

\subsection{The main result}

We are now equipped to describe the main construction of this section. We fix a metric $\phi\in\PSH(L)$, \ie a globally psh metric with logarithmic growth on $L$ and no identically $-\infty$ fibers (see Definition \ref{def_psh2}). Given any divisorial point $\nu_E$ associated to the component $E$ of a model $\cX$ of $X$, we know that $\phi+a\log|z|$ extends to a metric over $E$ for some $a\in\bbr$. Pick a psh metric $\phi_{E}$ with divisorial singularities of type $E$ on $\cX$, \ie locally of the form
\[
\phi_{E}=\log |f_{E}| + O(1),
\]
where $f_{E}$ is a local equation for $E$. We can then define a generic (signed) Lelong number
\begin{equation}\label{eq_genericlelong}
\varphi\na(\nu_E)=\ord_E(\phi):=-\sup\{c\geq 0,\,\phi+a\log|z|\leq c\cdot \phi_{E}+O(1)\text{ near E}\}+a.
\end{equation}
By linearity, this is independent of the choice of such an $a$. Performing this construction over all possible $E$ captures the singularities of $\phi$ along all possible models of $\cX$. Our main result for this section (which, as explained in the introduction, can be seen as parallel to some results of \cite{bfjlocal}) is then the following:
\begin{theorem}\label{thm_namap}
Let $X$ be a degeneration together with a relatively ample line bundle~$L$. The Lelong numbers of a metric $\phi\in\PSH(L)$ define a function on $X^{\mathrm{div}}$, which admits a unique $L$-psh extension, giving a map
\[
(\cdot)\na:\PSH(L)\to \PSH(L_\K\an),
\]
which is furthermore lower semicontinuous and order-preserving.
\end{theorem}

\subsection{Some preliminaries}\label{subsect_23}

We now prove some auxiliary results that will be useful in the proof of Theorem \ref{thm_namap}. We first show that multiplier ideals of psh metrics on~$L$ give $L_\K\an$-psh functions.

\begin{lemma}\label{lem_multidealisvertical}
Let $\phi$ be a metric in $\cE^1(L)$. Let $(\cX,\cL)$ be a model of $(X,L)$ such that $\phi$ extends as a psh metric on $\cL$. Then, up to restricting to a slightly smaller disc, for all $m$, the multiplier ideal $\ideal_m=\cJ(m\phi)$ is vertical, and there exists an integer $m_0$ (depending only on $\cL$ and not on $m$ or $\phi$) such that $(m+m_0)\cL\otimes \cJ(m\phi)$ is globally generated on $\cX$.
\end{lemma}
\begin{proof}
Since $\phi$ has fiberwise finite energy, it has zero Lelong numbers on all fibers, hence on $\cX-\cX_0$, as Lelong numbers cannot increase upon evaluating them on a larger space. Skoda's integrability theorem (\cite[Lem.\,5.6(a)]{demanalytic}) then yields, for all positive integers $m$, $L^1$-integrability of $e^{-m\phi}$. By \cite[Lem.\,5.6(a)]{demanalytic} again, the multiplier ideals satisfy $\ideal_m{}_{,x}=\cO_\cX{}_{,x}$ for all $m$ and for all $x$ outside of the central fiber, \ie $\ideal_m$ is cosupported on the central fiber.

Now, the global generation statement, follows from a relative equivalent of \cite[Prop.\,6.27]{demanalytic}. We can in fact argue just as in \cite[Lem.\,5.6]{bbj}: we must prove that there exists $m_0$ such that the sheaf $(m+m_0)\cL\otimes \cJ(\phi)$ is $\pi$-globally generated. By the relative Castelnuovo-Mumford criterion, having picked a relatively very ample line bundle $V$ on $X$ and an $m_0$ such that $m_0\cdot \cL - K_\cX - (d+1)V$ is relatively ample (after possibly restricting to a smaller disc), it is enough to show that for all $j=1,\dots,d$,
\[
R^j\pi_*(((m+m_0)\cL-jV)\otimes \cJ(m\phi))=0
\]
on the disc, which follows from Kodaira and Nadel vanishing.
\end{proof}

We thus obtain the following:
\begin{corollary}\label{cor_multidealvertical}For any metric $\phi\in\PSH(L)$, and any model $(\cX,\cL)$ of $(X,L)$ such that $\phi$ extends as a psh metric on $\cL$, there exists an integer $m_0$ such that the function
\[
(m+m_0)\mi\log |m\cJ(\phi)|
\]
is $\cL$-psh for all positive integers $m$.
\end{corollary}
\begin{proof}In the case where $\phi$ also has fiberwise finite energy, this follows from the previous lemma. In the general case, one can approximate $\phi$ on $\cL$ by a decreasing sequence of locally bounded metrics $\phi_k$. Since the integer $m_0$ depends only on $\cL$, the sequence
\[
k\mto \tilde\phi_k:=(m+m_0)\mi\log |m\cJ(\phi_k)|
\]
is then a sequence of $\cL$-psh functions. Since the sequence $\phi_k$ is decreasing, we have for all $k$ that $\cJ(\phi_{k+1})\subseteq \cJ(\phi_k)$, \ie the sequence $\tilde\phi_k$ is also decreasing, which implies its limit $(m+m_0)\mi\log |m\cJ(\phi)|$ to be $\cL$-psh, as desired.
\end{proof}

We conclude our preliminaries by introducing the log discrepancy function. Given a model $\cX$ of $X$, let $\cX^{\mathrm{div}}$ be the set of \textit{vertical} divisorial valuations on $\cX$, \ie of~divisorial valuations of the form $\nu_E=\ord_E$ with $E$ a divisor in the central fiber of a model $\cX'$ dominating $\cX$. Then, we can define its \textit{log discrepancy}
\[
A_\cX:\cX^{\mathrm{div}}\to \bbr,
\]
which will be used in the proof of Theorem \ref{thm_namap}. Let $\rho:\cY\to\cX$ be some model dominating $\cX$. The log discrepancy function of $\cX$ is then fully characterized by the formula
\[
K_{\cY}+\cY_0 = \rho^*(K_\cX + \cX_0) + \sum_i A_\cX(\nu_{E_i})a_i E_i,
\]
where $\cY_0=\sum_i a_i E_i$, and $K_\cX$, $K_\cY$ are the canonical bundles of $\cX$ and $\cY$ respectively. Since all the points in $\cX^{\mathrm{div}}$ arise in this way, this defines the log discrepancy function on $\cX^{\mathrm{div}}$ (see also \cite[\S 5.4]{bjtrop}).

\subsection{Proof of Theorem \ref{thm_namap}}\label{subsect_24}
We fix a metric $\phi\in\PSH(L)$. We need to show that the function defined on $X^{\mathrm{div}}$ by\vspace*{-3pt}
\[
\phi\na:\nu_E\mto \ord_E(\phi),
\]
where $\nu_E$ corresponds to a divisorial valuation and $\ord_E$ is defined as a generic Lelong number as in \eqref{eq_genericlelong}, admits a psh extension to $X_\K\an$. As a non-Archimedean psh function is uniquely defined on the set of divisorial points, it is then enough to show that $\phi\na$ can be approximated by a decreasing sequence of psh model functions on $X_\K\an$. Note that, by construction, the map $\phi\mto \phi\na$ is lsc and order preserving. By~Corol\-lary~\ref{cor_multidealvertical}, the metric\vspace*{-3pt}
\[
\psi_m=(m+m_0)\mi u_m,
\]
where $u_m$ is the model function $\log|\cJ(m\phi)|$, is $\cL_\R\an$-psh. Pick a divisorial point $\nu_E\in X^{\mathrm{div}}$ associated to a component in the central fiber of an analytic model $(\cX,\cL)$ of $(X,L)$. Now, one shows that\vspace*{-3pt}
\begin{equation}\label{eq_thatestimate}
m\cdot \varphi\na(\nu_E) \leq u_m(\nu_E) \leq m\cdot \varphi\na(\nu_E) + A_\cX(\nu_E),
\end{equation}
where $A_\cX$ is the log discrepancy function of $\cX$ as before, exactly as in \cite[Lem.\,5.7]{bbj}. Namely, by \cite[Lem.\,B.4]{bbj} applied to $\cX$, we have that, for any divisorial valuation $\nu_E$ in $\cX$,\vspace*{-3pt}
\[
\nu_E(\cJ(m\phi))\leq \nu_E(\phi)\leq \nu_E(\cJ(m\phi))+A_\cX(\nu_E),
\]
from which \eqref{eq_thatestimate} follows by definition of $\phi\na$ and $u_m$.

Having established \eqref{eq_thatestimate}, we deduce that the sequence $\psi_m$ is a sequence of $\cL_\R\an$-psh functions converging pointwise on $X^{\mathrm{div}}$ to $\varphi\na$. To show that $\varphi\na$ is $\cL_\R\an$-psh, it is then enough to prove that we can have this sequence be decreasing. By subadditivity of multiplier ideals (the main theorem in \cite{del}) we have $\cJ(2m\phi)\subseteq \cJ(m\phi)^2,$
thus
\[
\psi_{2m}\leq 2\psi_{m},
\]
and as $\phi_{m}\leq 0$,\vspace*{-3pt}
\[
\psi_{2m}\leq \frac{2(m+m_0)}{2m+m_0}\,\psi_m\leq \psi_m.
\]
Picking the subsequence $i\mto \psi_{2^i}$ therefore yields a decreasing subsequence converging to $\varphi\na$, as desired. We then set
$\phi\na:=\varphi\na + \phi_\cL,$
which concludes our proof.\qed

\subsection{Non-Archimedean limit of locally bounded metrics}\label{subsect_25}

We now begin studying the behaviour under the map $(\cdot)\na$ of the class of metrics $\phi$, such that there exists a model $(\cX,\cL)$ of $(X,L)$ on which $\phi$ admits a locally bounded extension.

\begin{prop}\label{prop_extension_on_model}
Let $\phi\in\PSH(L)$. Then,
\begin{enumerate}
\item\label{prop_extension_on_model1} $\phi$ extends to a psh metric on a model $(\cY,\cM)$ of $(X,L)$ if and only if $\phi\na\leq \phi_\cM$;
\item\label{prop_extension_on_model2} if furthermore the metric $\phi$ extends to a locally bounded psh metric on $(\cY,\cM)$, then $\phi\na=\phi_\cM$.
\end{enumerate}
\end{prop}
\begin{proof}
Note that it is equivalent to show the following: given an analytic model $(\cX,\cL)$ of $(X,L)$ and $\psi$ be a reference metric admitting a locally bounded extension (in particular, smooth) to $\cL$,
\eqref{prop_extension_on_model1} holds if and only if $\phi\na-\psi\na\leq \phi_\cM - \phi_\cL$, and \eqref{prop_extension_on_model2} if we have equality. This will allow us to work at the level of functions and relatively to another model, which is easier.

Assume first $\phi$ to extend to a psh metric on $\cM$. Let $\cZ$ dominate both models via $\pi_\cX:\cZ\to\cX$ and $\pi_\cY:\cZ\to \cY$. We have
\[
\pi_\cY^*\cM=\pi_\cX^*\cL+D
\]
for a unique Cartier divisor $D$ supported in the special fiber $\cZ_0$. Since $\phi$ extends to a psh metric on $(\cY,\cM)$ if and only if it extends to a psh metric on any model dominating $(\cY,\cM)$, we may without loss of generality focus on $\cZ$. Picking a local equation $f_D$ for the divisor $D$ obtained as above, $\phi$ extends to $\pi_\cY^*\cM$ if
\[
\phi-\psi\leq -\log|f_D| + C
\]
near $\cZ_0$. Taking generic Lelong numbers with respect to the underlying divisor of a divisorial point $\nu$ gives
$\nu(\phi)-\nu(\psi) \geq -\nu(D),$
\ie
\[
\phi\na(\nu)-\psi\na(\nu)\leq \phi_\cM(x)-\phi_\cL(x).
\]
In the case where $\phi$ admits a locally bounded extension, then there is also a lower bound, which shows by the same argument that $\phi\na=\phi_\cM-\phi_\cL$. The converse is obtained by uniqueness of the Siu decomposition of $\phi$ on $\cX$.
\end{proof}

\section{Finite-energy spaces and the Monge-Ampère extension property}\label{sect_sect3}

\subsection{Non-Archimedean finite-energy metrics}\label{sect_naenergy}

We begin this section with some reminders from non-Archimedean pluripotential theory. Let $X$ be a general variety over $\K$ endowed with an ample line bundle $L$. As in complex geometry, one can define Monge-Ampère measures associated to a tuple of $d=\dim X$ continuous non-Archimedean $L$-plurisubharmonic metrics. The general construction relies on intersection pairings (see \eg \cite{gublerinventionestropical}, \cite{be20}), or Chambert-Loir and Ducros' theory of differential forms on Berkovich spaces (\cite[5, 6]{cld}), building on Lagerberg's theory of differential superforms. We only describe the main results below.

Given $d$ continuous psh metrics $\phi_1,\dots,\phi_d$ on $L\an$, we have a Radon probability measure
\[
\MA(\phi_1,\dots,\phi_d)=V\mi\cdot dd^c\phi_1\wedge\dots\wedge dd^c\phi_d\wedge\delta_X,
\]
where $V=(L^d)$. For ease of notation, we will also write
\[
\MA(\phi)=V\mi\cdot dd^c\phi\wedge\dots\wedge dd^c\phi\wedge\delta_X=V\mi\cdot (dd^c\phi)^d\wedge\delta_X.
\]
Just as in the complex case, one can define the space of finite-energy metrics $\cE^1(L\an)$, having extended the Monge-Ampère energy via decreasing limits again.

Using the results of \cite{reb20b}, one can also metrize $\cE^1(L\an)$ via setting
\[
d_1(\phi_0,\phi_1)=E(\phi_0)+E(\phi_1)-2E(P(\phi_0,\phi_1)),
\]
where $P(\phi_0,\phi_1)$ is the envelope
\[
P(\phi_0,\phi_1)=\sup\,\{\phi\in \PSH(L),\,\phi\leq\min(\phi_0,\phi_1)\}.
\]
This gives $\cE^1(L\an)$ a metric space structure which is furthermore geodesic, and which admits distinguished maximal geodesics characterized by the fact that the energy is affine along them. Finally, much as in the complex setting, one can use non-Archimedean Deligne pairings over a point (\cite{be20}) to realize the relative Monge-Ampère energy between two metrics in $\cE^1(L\an)$:
\[
E(\phi_0)-E(\phi_1)=\langle \phi_0^{d+1}\rangle - \langle \phi_1^{d+1}\rangle.
\]
Finally, much as in Section \ref{sect_deligne}, we note that we can extend the Deligne pairing construction over a point in the non-Archimedean case, to line bundles metrized by non-Archimedean finite-energy metrics.

\subsection{The Monge-Ampère energy in the non-Archimedean limit}\label{subsect_32}

In the trivially-valued setting, we have already seen that a $\bbs^1$-invariant metric in $\hat\cE^1(L)$ coincides with a finite-energy psh geodesic ray $t\mto \phi_t$. Two natural ``asymptotic'' energies arise:
\begin{enumerate}
\item the radial limit $\lim_t \sfrac{E(\phi_t)}{t}$;
\item the non-Archimedean energy of the non-Archimedean metric $\phi\na$ associated to $\phi$.
\end{enumerate}
In \cite{bbj}, it is established that if $\phi$ extends to a locally bounded metric on a test configuration, then those two quantities coincide. This is not the case in general, however. In this section, we generalize those results to our relatively maximal psh metrics on degenerations. It will be clearer to express this using the relative dimension zero case of the construction from the previous section.

\begin{remark}[Relative dimension zero and the non-Archimedean limit]\label{rem_namap}As mentioned in Example \ref{exa_reldim2} and Corollary \ref{coro_trivialization}, given a model $(\cX,\cL)$ of $(X,L)$ and a metric $\phi\in\hat\cE^1(L)$, one can identify the Monge-Ampère energy $\langle \phi^{d+1}\rangle_{X/\pdisc}$ of~$\phi$ with a function on the punctured disc, by picking a trivialization $\tau$ of $\langle \cL^{d+1}\rangle$ and setting $u=-\log|\tau|_{\phi}$. The function $u$ then has a finite generalized slope (or Lelong number) at zero, but this Lelong number depends on the choice of a trivialization. A nice way of capturing all possible such Lelong numbers is by looking directly at the metric $(\langle \phi^{d+1}\rangle_{X/\pdisc})\na$ on $\langle L_\K^{d+1}\rangle$. The Lelong number of $u$ specifically is then recovered as the difference of Deligne pairings $(\langle \phi^{d+1}\rangle_{X/\pdisc})\na-\langle \phi_\cL^{d+1}\rangle$, where $\phi_\cL$ is the model metric associated to $\cL$ on $L_\K\an$.
\end{remark}

\begin{theorem}\label{thm_delignelelong}
For all $\phi\in\cE^1(L)$ admitting a locally bounded extension to some model $(\cX,\cL)$, we have
\[
(\langle \phi^{d+1}\rangle_{X/\pdisc})\na=\langle(\phi\na)^{d+1}\rangle,
\]
as non-Archimedean metrics on the Deligne pairing $\langle L\an \rangle$ over $\Spec \K$.
\end{theorem}
\begin{proof}
Note that the metric $\langle \phi^{d+1}\rangle_{X/\pdisc}$ is subharmonic by \eqref{eq_ddc}, so that the left-hand side is well-defined (this is the relative dimension zero case of Example \ref{exa_reldim2}).

We pick a model $(\cX,\cL)$ such that $\phi$ extends to a locally bounded metric on $\cL$. By~Proposition \ref{prop_extension_on_model}, we necessarily have $\phi\na=\phi_\cL$, the model metric on $L_\K\an$ associated to $\cL$, so that we are left to show that, given a trivialization $\tau$ of $\langle \cL^{d+1}\rangle_{\cX/\disc}$ and setting $u(z)=-\log|\tau(z)|_{\langle \phi_z^{d+1}\rangle}$, we have
$\nu_0(u)=0$
(recall how we defined the model metric $\phi_\cL$ in Section \ref{sect_nappt}). But $\phi$ is locally bounded near the central fiber of $\cL$, so that $u$ is locally bounded near zero, which implies $\nu_0(u) = 0$ as desired.
\end{proof}

\begin{remark}\label{rem_functionals}
We will occasionally refer to a metric satisfying the statement of Theorem \ref{thm_delignelelong} as satisfying the Monge-Ampère extension property. We also remark that the proof of the theorem works more generally for arbitrary Deligne pairings: given $d+1$ pairs of relatively ample line bundles $L_i$ on $X$ and metrics $\phi_i\in \cE^1(L_i)$ admitting locally bounded extensions to some model of $L_i$, one has
\[
(\langle \phi_0,\dots,\phi_d\rangle_{X/\pdisc})\na=\langle\phi\na_0,\dots,\phi\na_d\rangle.
\]
The fact that the slopes are well-defined follows as in the proof of the above theorem from the general property (\ref{eq_ddc}) of Deligne pairings. In Section \ref{subsect_35}, we will show how to extend this result to the class of metrics satisfying the Monge-Ampère extension property.
\end{remark}

\subsection{Hybrid maximal metrics: existence and uniqueness}\label{subsect_33}

We now study hybrid maximal metrics. Such metrics can be described as being relatively maximal, but with boundary values prescribed both at the complex boundary of $X$ and at the ``asymptotic'' or non-Archimedean boundary. We will then see that they correspond exactly to metrics satisfying the Monge-Ampère extension property.
\begin{defi}Let $\phi\in\hat\cE^1(L)$. We say that $\phi$ is hybrid maximal if for any $\psi\in\cE^1(L)$ such that $\psi\na \leq \phi\na$ and $\limsup(\psi-\phi)\leq 0$ near the boundary of $X$, we~have $\psi\leq\phi$.
\end{defi}

\begin{remark}
We show how to relate our terminology with that of \cite{bbj}, which deals with special cases of our objects:
\begin{itemize}
\item a geodesic ray in \cite{bbj} is a relatively maximal $\bbc^*$-invariant (logarithmic growth) psh metric on a line bundle over a test configuration in our article;
\item a maximal geodesic ray in \cite{bbj} is a hybrid maximal $\bbc^*$-invariant (logarithmic growth) psh metric on a line bundle over a test configuration in our article.
\end{itemize}
The ``hybrid'' refers to \eg the work of Boucksom-Jonsson, in which a hybrid property is a property that passes well from the complex setting to the non-Archimedean limit. Other possible denominations could be ``Lelong-maximal'' or ``maximal in the non-Archimedean limit'', but both of those seem to focus more on the limit behaviour while we require our metric to also be maximal in the complex world.
\end{remark}

Recall that, for the definition of $\cE^1(L)$, we had fixed some boundary data $\phi_\partial$, which is the restriction to $\pi\mi(\bbs^1)$ of a smooth psh metric on $X$.
\begin{theorem}\label{thm_uniquenesstropmaxdisc}
For any $\Phi\na\in \cE^1(L\an)$, there exists a unique metric $\phi\in\hat\cE^1(L)$ such that $\phi\na=\Phi\na$. In general, the space $\cE^1(L)$ is mapped by $(\cdot)\na$ to $\cE^1(L_\K\an)$; and, for any $\phi\in\cE^1(L)$, we have
\[
(\langle \phi^{d+1}\rangle_{X/\pdisc})\na\leq \langle (\phi\na)^{d+1}\rangle.
\]
\end{theorem}
\begin{remark}In other words, we do not have non-Archimedean extension of the Monge-Ampère energy in $\cE^1(L)$ in general, but simply an inequality.
\end{remark}
\begin{proof}[Proof. First step: the ample model case]
Assume $\Phi\na$ to be a model metric corresponding to an ample model $(\cX,\cL)$ of $(X,L)$. Define
\[
\phi:=\sup{}^*\{\psi\in\PSH(L),\,\lim_{z\to\xi}\psi(z)\leq \phi_{\partial}(\xi)\,\text{ for all }\xi\in \partial X,\,\psi\na\leq\Phi\na\}.
\]
Note that the class of metrics over which the supremum is taken is nonempty: since~$\cL$ is relatively ample, one can always pick a smooth psh metric $\psi$ on $\cL$, and by Proposition~\ref{prop_extension_on_model} the restriction of $\psi$ to $L$ belongs to this class (up to subtracting a large enough constant $c$ so that $\psi-c\leq \phi_\partial$ on the boundary).

We first show that $\phi\in\hat\cE^1(L)$. By Proposition \ref{prop_extension_on_model}, $\phi$ is the restriction to $L$ of the envelope
\[\textstyle
\tilde\phi:=\sup{}^*\{\psi\in\PSH(\cL),\,\lim_{z\to\xi}\psi(z)\leq \phi_{\partial}(\xi)\,\text{ for all }\xi\in \partial \cX\}.
\]
This envelope is (continuous) psh and relatively maximal by Theorem \ref{thm_existencemaximal}, and therefore so is $\phi$. Because $\phi$ is the supremum of metrics with logarithmic growth, it~also has logarithmic growth; and by continuity it naturally has fiberwise finite-energy. Thus, $\phi$ belongs to $\hat\cE^1(L)$. We now only have to argue hybrid maximality. This will follow from Proposition \ref{prop_extension_on_model}(1) and the extremal definition of $\phi$, if we can show there exists a metric $\psi\in\PSH(L)$ whose associated non-Archimedean metric~$\psi\na$ coincides with $\Phi\na$. But this is ensured by Proposition \ref{prop_extension_on_model}(1), since we can always construct a locally bounded metric on $\cL$. Thus $\phi$ satisfies the statement of the theorem.

\subsubsection*{Second step: an inequality for Lelong numbers of Monge-Ampère energies}
For the moment let $\phi\in\hat\cE^1(L)$, and let $\phi\na$ be its corresponding non-Archimedean metric. We approximate $\phi\na$ by a decreasing net of non-Archimedean model metrics $\phi_i\na$ corresponding to ample models $\cL_i$, to which are associated hybrid maximal metrics~$\phi_i$ on~$L$ by the first part of the proof. Because $\phi$ and the $\phi_i$ have logarithmic growth, we can by Lemma \ref{lem_extpshtame} pick a model $(\cX,\cL)$ of $(X,L)$ such that they extend quasi-plurisubharmonically to $\cL$. We then pick a trivialization $\tau$ of $\langle \cL^{d+1}\rangle_{\cX/\disc}$, and set
\[
E(\phi):z\mto E(\phi_z)=-\log|\tau(z)|_{\langle \phi_z^{d+1}\rangle},
\]
and define similarly $E(\phi_{i})$. By the first part of the proof,
\[
E(\phi_{i,z})=E\na(\phi_i\na)\log|z|+f(z)
\]
with $f(z)$ bounded on $\disc$ and depending only on the choice of boundary data for the definition of $\cE^1(L)$. On the other hand, because $\phi\na\leq\phi_i\na$, Proposition \ref{prop_extension_on_model} implies that $\phi$ defines a psh metric on $\cL_i$, hence there exists a constant $c$ with $\phi\leq \phi_i+c$. Therefore, much as in \cite[Lem.\,5.8]{bbj},
\[
E(\phi_{z})\leq E(\phi_{i,z})+c= E\na(\phi_i\na)\log|z|+f(z)+c,
\]
thus $\nu_0(E(\phi))\leq \nu_0(E(\phi_i))=E\na(\phi_i\na)$. We can as in Theorem \ref{thm_namap} up to a divisible enough subnet assume that $\phi_i\na$ is decreasing, so that by continuity of $E\na$ along decreasing limits, we obtain
\begin{equation}\label{eq_eq2}
\nu_0(E(\phi))\leq E\na(\phi\na)
\end{equation}
for all $\phi\in \hat\cE^1(L)$.

\subsubsection*{Third step: the general case}
We pick a net of model metrics $\phi_i\na$ associated to ample models, and decreasing to $\Phi\na$. Let $\phi_i\in\hat\cE^1(L)$ be the corresponding hybrid maximal metrics given by the first part of the proof. By maximality, the net $(\phi_i)_i$ is decreasing: let $\phi$ be the limit of this net. By Lemma \ref{lem_extpshtame}, $\PSH(L)$ is stable under decreasing limits, so that $\phi$ has maximal growth. Then, we can pick its associated non-Archimedean metric $\phi\na$. Note that it has finite energy: using the notations of the second part of the proof, for all $i$,
\[
E(\phi_{i,z})=E\na(\phi_i\na)\log|z|+f(z)\geq E\na(\Phi\na)+f(z),
\]
so that $E(\phi_z)>-\infty$ for all $z$ on taking decreasing limits; and because
\[
E(\phi_{z})=E\na(\Phi\na)\log|z|+f(z)
\]
it follows from Corollary \ref{coro_trivialization} that $\phi$ is a relatively maximal metric. Thus, $\phi\in \hat\cE^1(L)$. As follows from the above equality, $\nu_0(E(\phi))= E\na(\Phi\na)$, so that to show that $\phi$ is our desired metric, we only need to prove the equality $\phi\na=\Phi\na$.

Since the mappings $\phi\mto \phi\na$ are order-preserving, we find $\phi\na\leq\phi_k\na$
for~all~$k$,~\ie
\[
\phi\na\leq\Phi\na.
\]
By \cite[Prop.\,6.3.2]{reb20b}, we must then show that $E\na(\phi\na)=E\na(\Phi\na)$. By monotonicity of $E\na$ we already have
\[
E\na(\phi\na)\leq E\na(\Phi\na).
\]
Recalling that $\nu_0(E(\phi))=E\na(\Phi\na)$ it then suffices to show that
\[
\nu_0(E(\phi))\leq E\na(\phi\na).
\]
But this is given by \eqref{eq_eq2} and thus shows that $E(\phi\na)=E(\Phi\na)$, hence $\phi\na=\Phi\na$. This concludes the proof of the first statement, \ie the existence (and uniqueness) of a hybrid maximal metric giving the correct non-Archimedean metric.

\subsubsection*{Fourth step: the inequality in the case of a metric in $\cE^1(L)$}
We pick a metric $\phi\in\cE^1(L)$, and we define $\Phi$ to be the hybrid maximal metric with $\Phi\na=\phi\na$ obtained from the first statement of this theorem. Then, since $\Phi$ is relatively maximal, $\phi\leq \Phi$, so that by monotonicity of $\phi\mto\phi\na$ (and again keeping the energy notation from the second step of the proof),
\[
(E(\phi))\na\leq (E(\Phi))\na,
\]
while $(E(\Phi))\na=E\na(\Phi\na)$ by hybrid maximality, proving our inequality. Note that the logarithmic growth condition built into $\cE^1(L)$ forces $(E(\phi))\na$ to be finite, ensuring that $\phi\na$ belongs to $\cE^1(L_\K\an)$.
\end{proof}

\subsection{The isometric embedding}\label{subsect_34}

We denote by
\[
\hat\cE^1_{\mathrm{hyb}}(L)
\] the subspace of hybrid maximal metrics in $\hat\cE^1(L)$. Again we fix the following boundary condition for metrics in $\hat\cE^1(L)$: they are equal to $\phi_\delta$ on $\pi\mi(\bbs^1)$, where $\phi_\delta$ is the restriction of a smooth metric on $L$. Our main theorem is the following:
\begin{theorem}\label{thm_isometry}
The inverse of the mapping $(\cdot)\mto(\cdot)\na$ given by Theorem \ref{thm_uniquenesstropmaxdisc} is an isometric embedding of $(\cE^1(L\an),d_1\na)$ into $(\hat\cE^1(L),\hat d_1)$ with image $\hat\cE^1_{\mathrm{hyb}}(L)$. In~particular, $\hat\cE^1_{\mathrm{hyb}}(L)$ and $(\cE^1(L\an),d_1\na)$ are complete, geodesic metric spaces.
\end{theorem}
\begin{remark} The first statement of the theorem can be thought of as saying that hybrid maximal metrics have the $d_1$-extension property. In general, Theorem~\ref{thm_isometry} essentially means that we realize the (non-Archimedean) space $\cE^1(L\an)$ as a purely complex geometric object.
\end{remark}

\begin{proof}
Assume first the isometry statement to hold. Because $(\cE^1(L\an),d_1\na)$ is closed and the map is an isometry, $\hat\cE^1_{\mathrm{hyb}}(L)$ is a $\hat d_1$-closed subspace of $\hat\cE^1(L)$. On the other hand, the latter is $\hat d_1$-complete by Theorem \ref{thm_1completeness}, which implies $\hat\cE^1_{\mathrm{hyb}}(L)$ to be complete as well. Because the mapping is an isometry again, $(\cE^1(L\an),d_1\na)$ is then itself complete. Similarly, $\hat\cE^1_{\mathrm{hyb}}(L)$ is geodesic because the mapping is an isometry and $(\cE^1(L\an),d_1\na)$ is geodesic by \cite[Th.\,A]{reb20b}.

We now prove the isometry statement. Pick $\phi_0$, $\phi_1$ in $\hat\cE^1_{\mathrm{hyb}}(L)$. We assume both metrics to be continuous, and the general result will proceed as usual from regularization. Using Theorem \ref{thm_uniquenesstropmaxdisc} together with the expressions of the distances and additivity of Lelong numbers,
\begin{align*}
d_1\na(\phi_0\na,\phi_1\na)&=\langle(\phi_0\na)^{d+1}\rangle+\langle(\phi_1\na)^{d+1}\rangle-2\langle P(\phi_0\na,\phi_1\na)^{d+1}\rangle,
\\
d_1(\phi_{0,z},\phi_{1,z})&=\langle\phi_{0,z}^{d+1}\rangle+\langle\phi_{1,z}^{d+1}\rangle-2\langle P(\phi_{0,z},\phi_{1,z})^{d+1}\rangle,
\end{align*}
we only have to show that
\begin{equation}\label{eq_eqntoprove}
(-\langle P(\phi_{0},\phi_{1})^{d+1}\rangle_{X/\pdisc})\na=-\langle P(\phi_0\na,\phi_1\na)^{d+1}\rangle.
\end{equation}
Recall that we have seen $z\mto \langle P(\phi_{0,z},\phi_{1,z})^{d+1}\rangle$ to be superharmonic, so that the left-hand side is well-defined, being a Lelong number of a subharmonic function. To prove~\eqref{eq_eqntoprove} we construct another relatively maximal metric as follows. Set some \hbox{$r\in (0,1)$}. We consider the relatively maximal metric $\psi_r$ on the preimage $U_r$ of the annulus $\{r\leq |z| \leq 1\}$ with boundary data given by $z\mto P(\phi_{0,z},\phi_{1,z})$ for $z\in\partial U_r$, which exists by Theorem \ref{thm_existencemaximal}. Note that for all $z\in \{r\leq z \leq 1\}$ we have
\[
\psi_{r,z}\leq P(\phi_{0,z},\phi_{1,z}).
\]
Indeed, by relative maximality, $P(\phi,\psi)\!\leq\! \phi,\psi$ on $\partial U_r$ implies that for all~\hbox{$z\!\in\! \{r\!\leq\! z \!\leq\! 1\}$}, $\psi_{r,z}\leq \phi_{r,z},\psi_{r,z}$ holds. Because $\phi_{r,z},\psi_{r,z}$ are psh, the inequality then follows from the definition of the envelope $P(\phi_{0,z},\phi_{1,z})$.

This inequality implies again by relative maximality that for $z\in\pdisc$, the sequence $r\mto \psi_{r,z}$, $r\leq |z|$, is decreasing as $r$ decreases. Indeed let $r'<r\in (0,1)$ and $z$ with $|z|=r$. We then have $\psi_{r',z}\leq P(\phi_{0,z},\phi_{1,z})$, which means that the restriction of $\psi_{r'}$ to $U_r\subset U_{r'}$ is a relatively maximal metric with boundary data smaller than that of $\psi_r$ (which equals exactly $P(\phi_{0,z},\phi_{1,z})$ at $z$), and the comparison principle implies $\psi_{r',z}\leq \psi_{r,z}$ on $U_r$. Therefore, the limit $\lim_{r\to 0}\psi_r=:\psi$ is still a relatively maximal metric. We first prove that
\begin{equation}\label{eq_eqn3}
P(\phi_0\na,\phi_1\na)= \psi\na.
\end{equation}
To that end, we first claim that $\psi$ realizes the supremum
\begin{equation}\label{eq_eqn4}
\psi=\sup\{\varphi\in\PSH(L),\,\varphi\leq\phi_0,\phi_1\}.
\end{equation}
Since $\psi$ is itself such a metric, it is enough to show that for all candidates $\varphi$, we have $\varphi\leq\psi$. But for all $z\in X$, since $\varphi_z\leq \phi_{0,z},\phi_{1,z}$, we have
$\varphi_z\leq P(\phi_{0,z},\phi_{1,z})$,
hence
$\varphi_z\leq \psi_{r,z}$
and finally
\[
\varphi_z\leq \lim_r \psi_{r,z}=\psi_z.
\]
We now conclude the proof of \eqref{eq_eqn3}, using throughout Theorem \ref{thm_delignelelong}. By the extremal characterization \eqref{eq_eqn4} of $\psi$, we have that $\varphi\na\leq \psi\na$ for all $\varphi\leq \phi_0,\phi_1$. In particular, since the construction is order-preserving, the hybrid maximal metric $\Psi$ with $\Psi\na=P(\phi_0\na,\phi_1\na)$ satisfies $\Psi\leq \psi$, so that
\[
P(\phi_0\na,\phi_1\na)\leq \psi\na,
\]
while on the other hand, $\psi\leq \phi_0,\phi_1$, hence $\psi\na\leq \phi_0\na,\phi_1\na$ and finally $\psi\na\leq P(\phi_0\na,\phi_1\na)$, establishing \eqref{eq_eqn3}.
We now show that
\begin{equation}\label{eq_finaleqn}
-(\langle \psi^{d+1}\rangle_{X/\pdisc})\na=(-\langle P(\phi_{0},\phi_{1})^{d+1}\rangle_{X/\pdisc})\na.
\end{equation}
We will slightly abuse notation and consider each Deligne pairing as a function on~$\bbd^*$ (or a subset thereof), by which we mean that we have fixed a model $(\cX,\cL)$ and subtract Deligne pairings of the associated model metric. By definition, the function
\[
f_r(z):z\mto \langle \psi_{r,z}^{d+1} \rangle_{X/\{z\}}
\]
has as boundary values $\langle \phi_\delta^{d+1} \rangle$ on $\bbs^1$ and $\langle P(\phi_{0},\phi_1)^{d+1}\rangle$ on $r\bbs^1$. In particular, for all $r>0$ we have
\[
\langle -P(\phi_{0},\phi_{1})^{d+1}\rangle_{X/\pdisc}\rangle\na=-\lim_{r\to 0}\frac{\int_{r\bbs^1}f_r(z)\,dz}{\log r},
\]
so that it suffices to show that
\begin{equation}\label{eq_eqn}
\lim_{r\to 0}\frac{\int_{r\bbs^1}f_r(z)\,dz}{\log r}=(\langle \psi^{d+1}\rangle_{X/\pdisc})\na.
\end{equation}
By maximality, on the annulus $\{r\leq |z|\leq 1\}$, the harmonic function $f:z\mto \langle \psi^{d+1}_z\rangle_{X/\{z\}}$ is a (uniform) limit as $r>r'\to 0$ of the harmonic functions $f_{r'}$. It~is a classical property of harmonic functions that their average functions are linear in~$\log r$, \ie of the form
\[
\int_{r\bbs^1}f(z)\,dz=a_0\log r+b;\,\int_{r\bbs^1}f_{r'}(z)\,dz=a_{r'}\log r+b,
\]
from which one easily deduces that $b=\int_{\bbs^1} \langle \phi_\delta(z)^{d+1}\rangle_{X/\{z\}}\,dz$,
\[
a_0=(\langle \psi^{d+1}\rangle_{X/\pdisc})\na;\,a_{r'}=\int_{r'\bbs^1}f_{r'}(z)\,dz.
\]
Because $f_{r'}\to f$ uniformly on each annulus $\{r\leq |z|\leq 1\}$ for $r'<r$ we have $a_{r'}\to a_0$. which proves \eqref{eq_eqn}, hence \eqref{eq_finaleqn}. Equation \eqref{eq_finaleqn} together with \eqref{eq_eqn3} shows that \eqref{eq_eqntoprove} holds, concluding the proof.
\end{proof}

\begin{remark}The proof of the above result in the case of geodesic rays, which does not appear explicitly in the literature (but is based on some ideas from \cite{bdl}), was nicely explained to the author by Tamás Darvas.
\end{remark}

\begin{remark}In the above proof, we implicitly defined an envelope operator sending two metrics $\phi_0$, $\phi_1$ in $\hat\cE^1_{\mathrm{hyb}}(L)$ to the largest metric $\hat P(\phi_0,\phi_1)$ in $\hat\cE^1_{\mathrm{hyb}}(L)$ bounded above by $\phi_0$ and $\phi_1$. In \cite[Ex.\,3.3]{xia}, this construction appears already in the case of geodesic rays, and Xia uses this envelope to define the alternative distance
\[
\hat d_1'(\phi_0,\phi_1):=\lim_t t\mi (E(\phi_{0,t})+E(\phi_{1,t})-2E(\hat P(\phi_0,\phi_1)_t),
\]
which (a specialization of) our proof shows to coincide with the usual distance $\hat d_1$. In fact, Xia defines this envelope more generally in \cite[Ex.\,3.2]{xia}, in the radial equivalent of the space $\hat\cE^1(L)$. It is likely that this construction generalizes to metrics in $\hat\cE^1(L)$ in our setting, although this is outside the scope of the present article.
\end{remark}

\subsection{Non-Archimedean extension of generalized functionals}\label{subsect_35}

In Theorem \ref{thm_eonepairing}, we have seen that the fiberwise finite energy condition is the adequate condition for finiteness of fiberwise Deligne pairings. Further following the mantra that properties pertaining to the energy govern the same properties for more general Deligne pairings, we show that non-Archimedean extension of generalized energy functionals in the sense of Remark \ref{rem_functionals} holds for our class of hybrid maximal metrics, \ie metrics satisfying the Monge-Ampère extension property.

\begin{prop}\label{prop_functionals}Suppose given $d+1$ relatively ample line bundles $L_i$ on $X$. Then, for any $(d+1)$-uple of metrics $\phi_i\in\hat\cE^1_{\mathrm{hyb}}(L_i)$, we have
\[
(\langle\phi_0,\dots,\phi_d\rangle_{X/\pdisc})\na=\langle\phi\na_0,\dots,\phi\na_d\rangle.
\]
\end{prop}
\begin{proof}
We approximate each of the $\phi_i\na$ by a decreasing sequence of model metrics $\phi_{i,k}\na$, and denote by $\phi_{i,k}$ their associated hybrid maximal metrics. By our previous results, $\phi_{i,k}$ decreases to $\phi_i$ by hybrid maximality. The desired convergence of slopes now follows from some estimates from \cite[App.\,A]{bbj}, as in \cite[\S 4]{liytd}, which we sketch now. To that end we introduce the relative I-functional on $\cE^1(L_z)$:
\[
I(\phi,\psi):=\langle \phi-\psi,\psi^n\rangle_{X/\{z\}}-\langle \phi-\psi,\phi^n\rangle_{X/\{z\}}\geq 0.
\]
Now, the difference $\langle \phi_{0,k,z},\dots,\phi_{d,k,z}\rangle-\langle\phi_{0,z},\dots,\phi_{d,z}\rangle$ can be expressed, via the change of metric formula, as a sum of expressions of the form
\[
0\leq I_i:= \int_X (\phi_{i,k,z}-\phi_{i,z})(dd^c\phi_{0,z}\wedge\dots\wedge dd^c\phi_{i-1,z}\wedge dd^c \phi_{i+1,k,z}\wedge\dots\wedge dd^c\phi_{n,z}).
\]
We estimate each individual such term via \cite[Lem.\,A.2]{bbj}:
\begin{align*}
0\leq I_i\leq C\cdot d_1(\phi_{i,k,z},\phi_{i,z})^{2^{-n}}D(z),
\end{align*}
with
\[
D(z)=\max(I(\phi_{i,k,z}),I(\phi_{i,z}),\max_{j<i}(I(\phi_{j,z})),\max_{j>i}(I(\phi_{j,k,z})))^{1-2^{-n}},
\]
writing $I(\phi_{i,z}):=I(\phi_{i,z},\phi_{\mathrm{ref},z}^i)$ for some fixed reference metrics $\refmetric^i$ on each $L_i$. As in \cite[\S 4]{liytd}, one proceeds to see that the $I$-terms on the right-hand side are bounded by $C'\log|z|$ where $C'$ is a constant independent of $k$, while by the $d_1$-extension property of hybrid maximal metrics, the $d_1$-terms are bounded by $d_1\na(\phi_{i,k}\na,\phi_{i}\na)(-\log|z|)$, so~that
\[
0\leq \langle\phi_{0,k}\na,\dots,\phi_{d,k}\na\rangle-(\langle\phi_{0},\dots,\phi_{d}\rangle_{X/\pdisc})\na\leq C''\cdot \max_i d_1\na(\phi_{i,k}\na,\phi_{i}\na),
\]
and taking the limit in $k$ in the above inequality finally gives our result.
\end{proof}

\begin{example}\label{exa_functionals}Many functionals acting on $\PSH(L)$ satisfy the statement of the above Proposition. Having fixed some reference metric $\refmetric\in\hat\cE^1_{\mathrm{hyb}}(L)$, some among the most important are:
\begin{enumerate}
\item the $I$-functional, which appeared in the estimates mentioned in the above proof, which has many important norm-like properties and is commonly used to study properties of finite-energy spaces;
\item the $J$-functional, defined as
\[
J(\phi)=\langle\phi,\refmetric^d\rangle_{X/\pdisc}-\langle\refmetric^{d+1}\rangle_{X/\pdisc}-(d+1)\mi(E(\phi)-E(\refmetric)),
\]
which can be seen as a corrected relative Monge-Ampère energy which is translation invariant;
\item twisted energy functionals, defined as $E^{\psi}(\phi)=\langle \psi,\phi^d\rangle_{X/\pdisc},$
for $\psi\in\hat\cE^1_{\mathrm{hyb}}(L')$, where $L'$ is another line bundle on $X$. A special case of it appears in the expression of the Mabuchi K-energy, and the study of its slopes in the trivially-valued case is essential to establish the general (cscK) case of the Yau-Tian-Donaldson conjecture, as in \cite{liytd}.
\end{enumerate}
\end{example}

\subsection{Test configurations and the trivially valued case}\label{subsect_36}

All of our previous results encapsulate the trivially valued case, as we explain now. Let $\pi:X\to\bbd^*$ be now a polarized test configuration, \ie a degeneration with relatively ample line bundle~$L$ such that $\pi$ and $L$ are equivariant under some $\bbc^*$-action (forcing all fiber pairs $(X_z,L_z)$ to be isomorphic). One may then choose a reference continuous psh metric $\refmetric$ on the fiber at $1$ and require our psh metrics $\phi$ to satisfy $\phi_z=i_z{}^*\refmetric$ for $z\in\bbs^1$, and with
\[
i_z:X_z\to X_1
\]
the isomorphism as mentioned above. The authors in \cite{bbj} study the space $\cE^1_0(X_1\an)$ of finite-energy metrics over the analytification of $X_1$ with respect to the trivial absolute value on $\bbc$. We denote by $\cR^1(L_1)$ the space of hybrid maximal finite-energy rays in $\PSH(L_1)$ emanating from $\refmetric$ (where, as mentioned before, a hybrid maximal ray corresponds in the terminology of \cite{bbj} to a maximal psh geodesic ray). We then claim the following:
\begin{prop}
There is a sequence of distance-preserving maps
\[
\cE^1_0(L_1\an)\simeq \cR^1(L_1) \hto \hat\cE_{\mathrm{hyb}}^1(L) \simeq \cE^1(L\an),
\]
where the first and last maps are bijective (\ie isometries), while the middle map is injective.
\end{prop}
\begin{proof}
The case of the last map has been treated by Theorem \ref{thm_isometry}. The rest of the proof is merely a matter of correctly defining our maps.

For the first map, the bijection is given by \cite[Th.\,6.6]{bbj}. The metrization of the space $\cE^1_0(L_1\an)$ is described in a \cite{bjtrivval21}, but proceeds much as the metrization of $\cE^1(L\an)$ in \cite{reb20b}, while we recall that we metrize the space of maximal rays by
\[
\hat d_{1,0}(\phi,\phi')=\lim_t t\mi d_1(\phi_t,\phi'_t)
\]
and take equivalence classes to yield the space $\cR^1(L_1)$.
(We direct the reader to \eg \cite{bdl}. Note that in the cited article, the authors consider the space of all (non-necessarily hybrid) maximal psh rays.) Proving the distance-preserving-ness of the isomorphism is then essentially a simpler version of Theorem \ref{thm_isometry}, which we leave to the interested reader.

We claim that the middle map, which we will denote $\iota_0$, can be represented as follows: let $\phi:t\mto \phi_t$ be a hybrid maximal psh geodesic ray in $X_1$. Let $i_z$ be as before the isomorphism $i_z:X_z\to X_1$, and define $\iota_0(\phi)$ to be the metric
$z\mto i_z{}^*(\phi_{-\log|z|}).$
The distance-preservingness is immediate, so that we are left to check that $\iota_0(\phi)$ is a hybrid maximal metric. By \cite[Cor.\,6.7]{bbj},
$t\mto E(\phi_t)$
is affine, which implies by invariance of the energy under polarized isomorphisms that
$z\mto \iota_0(\phi)(z)$
is harmonic on $\bbd^*$, proving maximality by Proposition \ref{prop_harmonicitymaximaldiscs}, and hybrid maximality is given by construction.
\end{proof}

\subsection{Kähler-Einstein metrics in families}\label{subsect_misc}

Let $M$ be a complex projective manifold with ample canonical bundle $K_M$. It is a consequence of the by now classical Aubin-Yau theorem that $M$ carries a Kähler-Einstein metric. If $X$ is more generally a family of canonically polarized projective manifolds, the family $z\mto\phi_z$ of fiberwise Kähler-Einstein metrics is known to have plurisubharmonic variation (from the work of \eg Schumacher \cite{sch08}) and to have logarithmic growth (\cite[Th.\,3]{sch08}). In particular, $\phi$ defines a metric in our class $\cE^1(K_X)$. Interpreting the work of Pille-Schneider through our lens, \cite[Th.\,A]{ps} should conjecturally imply that $\phi\na$ corresponds to a model $K_X\an$-function associated to a distinguished model of $(X,K_X)$.

An immediate question arises: how does the metric $\phi$ relate to our relatively maximal metrics framework? In particular, how does $\phi$ relate to the hybrid maximal metric $\Phi$ corresponding to $\phi\na$? Interestingly, $\phi$ is not even relatively maximal when the Kodaira-Spencer class of the family $X$ is nontrivial, by \cite[Main\,Th.]{sch08}, since $\phi$ will be strictly positive (in particular, cannot satisfy $\MA(\phi)=0$). As a consequence, we have that $dd^c d_1(\phi,\Phi)=dd^c (E(\Phi) - E(\phi))$ is given explicitly by the formula of Schumacher, using the pushforward formula for Deligne pairings.

Naturally, it would be interesting to know whether one could detect via non-Archimedean tools the existence of a family of Kähler-Einstein metrics in the class of a hybrid maximal metric. This seems a bit ambitious, since one only captures the ``asymptotic'' behaviour of a family of metrics when considering non-Archimedean data. A more realistic (and perhaps just as interesting) problem would be solving the following hybrid ``almost Kähler-Einstein'' problem: to find $\phi\in\hat\cE^1_{\mathrm{hyb}}(K_X)$, such that $\phi\na=\psi\na$ where $\psi\na$ is an ``almost Kähler-Einstein metric'':
\begin{align*}
(\omega_z +i\partial\bar\partial \psi_z)^d - e^{h_{\omega_z} + \psi_z}\omega_z{}^d \to_{z\to 0} 0.
\end{align*}
The upshot is that this problem gives, intuitively, a purely non-Archimedean criterion for the existence of a family of complex manifolds degenerating to a Kähler-Einstein manifolds. (Of course, the same problem arises in the (possibly twisted) Fano case.)

Finally, we briefly mention an additional difficulty in the Calabi-Yau case. By a counterexample of Cao-Guenancia-Paun, we know that a family $\phi=(\phi_z)_z$ of Kähler-Einstein metrics on a degeneration of Calabi-Yau manifolds does not necessarily vary plurisubharmonically (\cite[Th.\,3.1]{cgp2}). One can however take the plurisubharmonic envelope $P(\phi)$ of $\phi$, and then the hybrid maximal metric $\Phi$ with $\Phi\na=P(\phi)\na$. In \cite{bjtrop}, Boucksom-Jonsson show that the family of measures $\MA(\phi_z)$ converge in a certain sense to the non-Archimedean Monge-Ampère measure of some metric $\psi\na$. We therefore formulate the following conjecture, which would connect our hybrid maximal setting with degenerations of Kähler-Einstein metrics on Calabi-Yau manifolds:
\begin{conjecture}
$\psi\na=P(\phi)\na=\Phi\na$.
\end{conjecture}

\newpage
\backmatter
\bibliographystyle{jepalpha+eid}
\bibliography{reboulet}
\end{document}
